\documentclass[10pt,letterpaper,openany]{book}

% ──────────────────────────────────────────────
% Packages
% ──────────────────────────────────────────────
\usepackage[margin=1.5in]{geometry}
\usepackage{amsmath, amssymb, amsthm}
\usepackage{mathtools}
\usepackage{bm}
\usepackage{tikz}
\usepackage{pgfplots}
\pgfplotsset{compat=1.18}
\usetikzlibrary{arrows.meta, decorations.markings, calc, patterns, positioning}
\usepackage{tcolorbox}
\tcbuselibrary{theorems, breakable, skins}
\usepackage{hyperref}
\usepackage{imakeidx}
\makeindex[intoc, title=Index]
\usepackage{enumitem}
\usepackage{fancyhdr}
\usepackage{titlesec}
\usepackage{epigraph}
\usepackage{xcolor}
\usepackage{booktabs}
\usepackage{listings}
\usepackage{../geometry_of_motion}

% ──────────────────────────────────────────────
% Code Formatting
% ──────────────────────────────────────────────
\definecolor{codegreen}{rgb}{0,0.6,0}
\definecolor{codegray}{rgb}{0.5,0.5,0.5}
\definecolor{codepurple}{rgb}{0.58,0,0.82}
\definecolor{backcolour}{rgb}{0.95,0.95,0.92}

\lstdefinestyle{mystyle}{
    backgroundcolor=\color{backcolour},
    commentstyle=\color{codegreen},
    keywordstyle=\color{magenta},
    numberstyle=\tiny\color{codegray},
    stringstyle=\color{codepurple},
    basicstyle=\ttfamily\footnotesize,
    breakatwhitespace=false,
    breaklines=true,
    captionpos=b,
    keepspaces=true,
    numbers=left,
    numbersep=5pt,
    showspaces=false,
    showstringspaces=false,
    showtabs=false,
    tabsize=2
}
\lstset{style=mystyle}
\gomapplylistingstyle

% ──────────────────────────────────────────────
% Colors
% ──────────────────────────────────────────────
\definecolor{chapblue}{RGB}{0, 70, 130}
\definecolor{accentred}{RGB}{180, 40, 40}
\definecolor{softgray}{RGB}{245, 245, 245}
\definecolor{trajgreen}{RGB}{30, 130, 60}
\definecolor{darkgold}{RGB}{160, 120, 20}
\definecolor{neuralblue}{RGB}{40, 80, 160}

% ──────────────────────────────────────────────
% Theorem environments
% ──────────────────────────────────────────────
\newtcbtheorem[number within=chapter]{neurobox}{Neural Architecture}{
  colback=neuralblue!5, colframe=neuralblue!80,
  fonttitle=\bfseries, separator sign={.~},
  breakable
}{neuro}

\newtcbtheorem[number within=chapter]{insight}{Key Insight}{
  colback=darkgold!5, colframe=darkgold!80,
  fonttitle=\bfseries, separator sign={.~},
  breakable
}{insight}

\newtcbtheorem[number within=chapter]{example}{Example}{
  colback=softgray, colframe=black!40,
  fonttitle=\bfseries, separator sign={.~},
  breakable
}{ex}

\newtcbtheorem[number within=chapter]{definition}{Definition}{
  colback=chapblue!5, colframe=chapblue!75,
  fonttitle=\bfseries, separator sign={.~},
  breakable
}{def}







\newtcbtheorem[number within=chapter]{laymansbox}{In Plain Language}{
  colback=green!5!white, colframe=green!50!black,
  fonttitle=\bfseries, separator sign={.~},
  breakable
}{laymans}

\makeatletter
\theoremstyle{plain}
\@ifundefined{theorem}{\newtheorem{theorem}{Theorem}[chapter]}{}
\@ifundefined{proposition}{\newtheorem{proposition}[theorem]{Proposition}}{}
\makeatother

% ──────────────────────────────────────────────
% Notation Concordance — Volume IV
% ──────────────────────────────────────────────
% This volume follows Volume 0 conventions with additions for motor control:
%   Vectors:    bold italic via \bm{}     e.g. \state = \bm{x}
%   Reals:      \Reals  (alias \R also available via geometry_of_motion.sty)
%   Additional: \drift = f, \inputfield = g, \configvec = \bm{q}
%
% See geometry_of_motion.sty for the canonical fallback definitions.
% ──────────────────────────────────────────────

% ──────────────────────────────────────────────
% Custom commands
% ──────────────────────────────────────────────
\renewcommand{\state}{\bm{x}}
\renewcommand{\control}{\bm{u}}
\renewcommand{\traj}{\bm{\gamma}}
\newcommand{\drift}{f}
\newcommand{\inputfield}{g}
\newcommand{\configvec}{\bm{q}}
\newcommand{\configdot}{\dot{\bm{q}}}
\newcommand{\configddot}{\ddot{\bm{q}}}
\renewcommand{\Reals}{\mathbb{R}}
\renewcommand{\dd}{\mathrm{d}}
\renewcommand{\dt}{\,\dd t}

\DeclareMathOperator*{\argmin}{arg\,min}
\DeclareMathOperator*{\argmax}{arg\,max}

\renewcommand{\norm}[1]{\left\lVert #1 \right\rVert}

% ──────────────────────────────────────────────
% Chapter formatting
% ──────────────────────────────────────────────
\titleformat{\chapter}[display]
  {\normalfont\huge\bfseries\color{chapblue}}
  {\chaptertitlename\ \thechapter}{20pt}{\Huge}
\titlespacing*{\chapter}{0pt}{-20pt}{30pt}

% ──────────────────────────────────────────────
% Header/Footer
% ──────────────────────────────────────────────
\pagestyle{fancy}
\fancyhf{}
\fancyhead[LE,RO]{\thepage}
\fancyhead[RE]{\textit{Volume IV: Human Motor Control}}
\fancyhead[LO]{\textit{\nouppercase{\leftmark}}}
\renewcommand{\headrulewidth}{0.4pt}
\setlength{\headheight}{14.5pt}

% ──────────────────────────────────────────────
% Title Page
% ──────────────────────────────────────────────
\title{%
  {\Huge\bfseries\color{chapblue} Tangent-Space Methods for Nonlinear Control and Biomechanics}\\[0.4cm]
  {\Large Volume IV: Human Motor Control\\
  The Neural Architecture of Movement}\\[0.6cm]
  {\color{chapblue}\rule{\textwidth}{2pt}}
}
\date{}

\hypersetup{colorlinks=true, linkcolor=chapblue, filecolor=magenta, urlcolor=chapblue, citecolor=accentred}


\begin{document}

\maketitle
\thispagestyle{empty}
\cleardoublepage

\frontmatter
\chapter*{Preface}
\addcontentsline{toc}{chapter}{Preface}

\begin{gomdraftnotice}
This volume is under active development. Chapter depth and validation vary;
some sections remain preliminary. The biology and nonlinear-dynamics chapter
includes worked examples and exercises with independently checked mechanics.
That review does not certify every chapter, a neural mechanism, or a human
performance model. Source and validation records distinguish the completed
corrections from the remaining review.
\end{gomdraftnotice}

Volumes~0--III developed the mathematical and physical foundations for understanding
motion: differential geometry, nonlinear control, trajectory optimization, and
biomechanics. This volume asks a fundamentally different question: \emph{how does the
brain actually control the body?}

The golfer coordinates a body, a flexible implement and changing contacts using
actuators whose force capacity depends on activation and mechanical history.
The relevant model dimension, movement duration, precision and feedback latency
depend on the task and measurement; they cannot be reduced to one universal
count or timing budget.

This volume investigates how parallel processing, recurrent feedback,
prediction, learned policies and intrinsic mechanics can cooperate. Cortical
histological layers are not sequential computational layers, and a synaptic
delay divided into a movement duration does not establish network depth or
prove one neural architecture. A low-dimensional or shallow controller can be
a useful engineering construction without uniquely explaining human control.

The purpose is to connect explicit dynamical models with observations that
could support or challenge them. For golf and humanoid manipulation alike,
the questions concern available control authority, coupled stability, energy
and task accuracy under stated physical and informational constraints.

\tableofcontents
\chapter*{Nomenclature}
\addcontentsline{toc}{chapter}{Nomenclature}

\section*{General Conventions}
\begin{itemize}
    \item Scalars are lowercase or Greek letters (e.g., $m, t, \theta$).
    \item Vectors are bold lowercase (e.g., $\bx, \bu, \bq$).
    \item Matrices are bold uppercase (e.g., $\bA, \bB, \bM, \bJ$).
    \item Expected values and sets are denoted by blackboard bold or script (e.g., $\mathbb{E}, \mathcal{S}$).
\end{itemize}

\section*{State and Kinematics}
\begin{description}
    \item[$\bq \in \Reals^n$] Joint configuration vector (generalized coordinates).
    \item[$\dot{\bq} \in \Reals^n$] Joint velocity vector (generalized velocities).
    \item[$\ddot{\bq} \in \Reals^n$] Joint acceleration vector.
    \item[$\bx \in \Reals^{nx}$] System state vector; commonly $\bx = [\bq^T, \dot{\bq}^T]^T$ for mechanical systems.
    \item[$\bu \in \Reals^m$] Control input vector (e.g., joint torques, muscle activations).
    \item[$\btau \in \Reals^n$] Generalized forces/torques acting on the joints.
    \item[$t \in \Reals$] Time.
\end{description}

\section*{Inertia and Dynamics}
\begin{description}
    \item[$\bM(\bq)$] Mass (inertia) matrix, symmetric positive definite.
    \item[$\bC(\bq, \dot{\bq})$] Coriolis and centrifugal force matrix.
    \item[$\bg(\bq)$] Gravity vector.
    \item[$\bJ(\bq)$] Jacobian matrix relating joint velocities to task-space velocities ($\dot{\bp} = \bJ(\bq)\dot{\bq}$).
\end{description}

\section*{Control and Optimization}
\begin{description}
    \item[$\bA_k = \frac{\partial f}{\partial \bx}$] Local Jacobian matrix of the state dynamics.
    \item[$\bB_k = \frac{\partial f}{\partial \bu}$] Local Jacobian matrix of the control input.
    \item[$V(\bx)$ or $\mathcal{V}$] Value function or Lyapunov function.
    \item[$\bQ$] State penalty matrix in LQR/DDP.
    \item[$\bR$] Control penalty matrix in LQR/DDP.
    \item[$\bK$] Feedback gain matrix ($\delta\bu = -\bK\delta\bx$).
    \item[$\traj$] A trajectory or flow, mapping time and initial conditions to states.
\end{description}

\section*{Biomechanics and Motor Control}
\begin{description}
    \item[$a_i \in [0, 1]$] Muscle activation level.
    \item[$\ell_i^M$] Muscle fiber length.
    \item[$v_i^M$] Muscle fiber contraction velocity.
    \item[$\ell_i^{MT}$] Musculotendon unit length.
    \item[$\bm{F}_{\text{muscle}}$] Force generated by muscles.
    \item[$\bW$] Muscle synergy matrix (weights).
    \item[$\bc(t)$] Muscle synergy activation vector over time.
\end{description}

\section*{Geometry and Manifolds}
\begin{description}
    \item[$SE(3)$] Special Euclidean group of rigid body motions.
    \item[$SO(3)$] Special Orthogonal group of 3D rotations.
    \item[$\mathfrak{se}(3)$] Lie algebra of $SE(3)$, space of spatial twists (velocities).
    \item[$\mathcal{M}$] A smooth manifold.
    \item[$T_{\bx}\mathcal{M}$] Tangent space at point $\bx$.
\end{description}

\cleardoublepage

\mainmatter

\chapter{The Degrees-of-Freedom Problem}
\label{ch:dof-problem}

\epigraph{%
  The coordination of movement is the process of mastering redundant
  degrees of freedom of the moving organ, in other words its conversion
  to a controllable system.
}{--- Nikolai Bernstein, 1967}

\section{Bernstein's Original Formulation}

In 1935, Nikolai Bernstein, a Soviet physiologist, made a deceptively simple
observation: the number of degrees of freedom available for any motor task
\emph{vastly exceeds} the number required to define the task. This observation,
formalized in his 1967 monograph \cite{Bernstein1967}, is the founding problem
of motor control.

\begin{definition}{The Degrees-of-Freedom Problem}{dof-problem}
  Let a motor task be defined by $k$ task-space constraints (e.g., $k = 6$ for
  placing a hand at a specific position and orientation). Let the body have $n$
  joint degrees of freedom and $p > n$ muscles. The \textbf{degrees-of-freedom
  problem} is: how does the nervous system select a specific solution
  $(\bm{q}(t), \bm{u}(t)) \in \Reals^n \times \Reals^p$ from the infinite set
  of solutions consistent with the task constraints?

  The \textbf{kinematic redundancy} is $n - k$ (extra joint configurations).
  The \textbf{muscle redundancy} is $p - n$ (extra muscles per joint).
  The \textbf{total redundancy} is $(n - k) + (p - n) = p - k$.
\end{definition}

\subsection{The Numbers Are Staggering}

Consider the golf swing:
\begin{center}
\renewcommand{\arraystretch}{1.3}
\begin{tabular}{@{}lr@{}}
\toprule
\textbf{Quantity} & \textbf{Approximate Value} \\
\midrule
Skeletal DOF involved & $\sim$200 \\
Muscles contributing & $\sim$400 \\
Task constraints (clubhead at impact) & 6 \\
Kinematic redundancy & $\sim$194 \\
Muscle redundancy & $\sim$200 \\
Total redundancy & $\sim$394 \\
Duration of downswing & $\sim$0.25~s \\
Neural processing delay & $\sim$0.1$\,$--$\,$0.2~s \\
\bottomrule
\end{tabular}
\end{center}

The nervous system must select from $\sim$400-dimensional muscle activation patterns
to satisfy 6 constraints, while accounting for interaction torques, gravity, and
neural delays---\emph{in a quarter of a second}. This is the problem.

\begin{figure}[h]
\centering
\begin{tikzpicture}[scale=1.2]
  % Task space (low-dim)
  \draw[->, thick] (0, 0) -- (1, 0) node[align=center, anchor=north] {Task DOF = 6};
  \draw[->, thick] (0, 0) -- (0, 1) node[align=center, anchor=east] {Position + Orientation};

  % Mapping to configuration space
  \draw (2.5, 0.5) ellipse (0.6 and 1.5) node[align=center, anchor=center] {\footnotesize Config};

  % Null space visualization
  \draw (5, 0) ellipse (1.5 and 2);
  \draw (5, 0.1) node[align=center, anchor=center] {Null space};
  \draw (5, -0.8) node[align=center, anchor=center] {\footnotesize (task-irrelevant};
  \draw (5, -1.2) node[align=center, anchor=center] {\footnotesize variability)};

  % Arrow showing redundancy
  \draw[->, thick, dashed] (1.8, 0.5) -- (2.9, 0.5) node[align=center, anchor=south, midway] {\small redundant DOF};
  \draw[->, thick] (3.2, 0.5) -- (3.5, 0.5);

  % Label
  \draw (2.5, 2.5) node[align=center, anchor=center] {\textbf{Bernstein's DOF Problem}};
  \draw (2.5, 2.1) node[align=center, anchor=center] {\small Many solutions satisfy the task constraints};
\end{tikzpicture}
\caption{The degrees-of-freedom problem: redundancy in configuration space allows
task-irrelevant variability in the null space of the task Jacobian.}
\label{fig:ch1:dof-illustration}
\end{figure}

\section{From Problem to Resource: Motor Abundance}

Modern motor control theory has reframed Bernstein's ``problem'' as an
\emph{opportunity}. Redundancy is not a curse but a resource:

\begin{insight}{Redundancy as Motor Abundance}{abundance}
  The extra degrees of freedom provide:
  \begin{enumerate}
    \item \textbf{Flexibility}: the same task can be performed in many ways,
          allowing adaptation to constraints (fatigue, injury, obstacles)
    \item \textbf{Robustness}: perturbations can be absorbed without task failure
    \item \textbf{Exploration}: variability in execution allows the system to
          discover better solutions over time
    \item \textbf{Error correction}: task-irrelevant DOF can compensate for
          errors in task-relevant DOF
  \end{enumerate}
  We adopt the term \textbf{motor abundance} (Latash 2012) to emphasize that
  redundancy is a feature, not a bug.
\end{insight}

\section{The Uncontrolled Manifold Hypothesis}

The uncontrolled manifold (UCM) hypothesis \cite{ScholzSchoner1999} provides a
mathematical framework for understanding how the nervous system manages redundancy.

\begin{definition}{The Uncontrolled Manifold}{ucm}
  Let $\bm{q} \in \Reals^n$ be the joint configuration and $\bm{y} = h(\bm{q}) \in \Reals^k$
  be the task variable. The \textbf{uncontrolled manifold} at a nominal configuration
  $\bm{q}_0$ is the null space of the task Jacobian:
  \begin{equation}
    \text{UCM}(\bm{q}_0) = \left\{ \delta\bm{q} \in \Reals^n :
    J_h(\bm{q}_0) \delta\bm{q} = \bm{0} \right\},
    \label{eq:ch1:ucm-def}
  \end{equation}
  where $J_h = \partial h / \partial \bm{q}$ is the task Jacobian. Variability within
  the UCM does not affect the task variable. Variability perpendicular to the UCM
  (in the \textbf{orthogonal complement}) directly affects task performance.
\end{definition}

\begin{laymansbox}{Mathematical Reminder: The Geometry of the UCM}{ucm_geometry_reminder}
Recall our differential geometry from Volume 0: a manifold is a curved surface representing valid states. The \textbf{Uncontrolled Manifold} (UCM) is quite literally a geometric submanifold embedded inside the larger configuration space $\mathcal{Q}$. 

Because the task geometry is non-linear (e.g., the arc of a swinging robot arm), the UCM is curved. We calculate its flat \textbf{Tangent Space} mathematically using the null space of the Jacobian ($J_h \delta\bm{q} = \bm{0}$). As long as the nervous system's noise vectors point exclusively into this tangent space, the physical position of the hand will not budge from the target. The nervous system is actively shaping its noise ellipsoid to align with the tangent space of the target manifold!
\end{laymansbox}

The UCM hypothesis predicts that the nervous system:
\begin{enumerate}
  \item \textbf{Suppresses} variability in the orthogonal complement of the UCM
        (task-relevant variability)
  \item \textbf{Tolerates} or even \textbf{encourages} variability within the UCM
        (task-irrelevant variability)
\end{enumerate}

This has been experimentally confirmed across many motor tasks \cite{ScholzSchoner1999,
Todorov2002}.

\begin{lstlisting}[language=Python, caption={Uncontrolled manifold analysis.}]
import numpy as np
from typing import Tuple

def ucm_analysis(
    joint_angles: np.ndarray,
    task_jacobian_func,
    n_repetitions: int | None = None
) -> Tuple[float, float, float]:
    """Perform UCM variance decomposition.

    Parameters
    ----------
    joint_angles : np.ndarray, shape (n_reps, n_dof)
        Joint angle data from multiple repetitions of the same task.
    task_jacobian_func : callable
        q -> J(q), shape (k, n_dof), the task Jacobian.
    n_repetitions : int, optional
        Number of repetitions (inferred from data if not given).

    Returns
    -------
    Tuple[float, float, float]
        (variance_ucm, variance_ort, ratio)
        UCM variance, orthogonal variance, and their ratio.
    """
    n_reps, n_dof = joint_angles.shape
    q_mean = np.mean(joint_angles, axis=0)

    # Task Jacobian at the mean configuration
    J = task_jacobian_func(q_mean)
    k = J.shape[0]  # task dimension

    # Null space of J (UCM directions)
    U, S, Vt = np.linalg.svd(J, full_matrices=True)
    rank = np.sum(S > 1e-10)
    null_space = Vt[rank:].T  # Shape: (n_dof, n_dof - rank)

    # Orthogonal complement (range space)
    range_space = Vt[:rank].T  # Shape: (n_dof, rank)

    # Decompose variability
    deviations = joint_angles - q_mean  # (n_reps, n_dof)

    # Project onto UCM
    proj_ucm = deviations @ null_space  # (n_reps, n_dof - rank)
    var_ucm = np.mean(np.sum(proj_ucm**2, axis=1))

    # Project onto orthogonal complement
    proj_ort = deviations @ range_space  # (n_reps, rank)
    var_ort = np.mean(np.sum(proj_ort**2, axis=1))

    # Normalize by dimensionality
    var_ucm_per_dof = var_ucm / (n_dof - rank) if n_dof > rank else 0
    var_ort_per_dof = var_ort / rank if rank > 0 else 0

    ratio = var_ucm_per_dof / var_ort_per_dof if var_ort_per_dof > 0 else np.inf

    return var_ucm_per_dof, var_ort_per_dof, ratio


# Example: 3-joint planar arm reaching to a point
def planar_arm_jacobian(q: np.ndarray, L: np.ndarray = np.array([0.3, 0.25, 0.2])):
    """Task Jacobian for endpoint position of a 3-link planar arm."""
    n = len(q)
    J = np.zeros((2, n))
    for i in range(n):
        for j in range(i, n):
            angle_sum = np.sum(q[:j+1])
            J[0, i] -= L[j] * np.sin(angle_sum)
            J[1, i] += L[j] * np.cos(angle_sum)
    return J

# Generate synthetic reaching data (50 repetitions)
np.random.seed(42)
q_nominal = np.array([0.5, 1.2, -0.3])
n_reps = 50
# Motor noise: more variability in task-irrelevant directions
noise = np.random.randn(n_reps, 3) * 0.05
q_data = q_nominal + noise

v_ucm, v_ort, ratio = ucm_analysis(q_data, planar_arm_jacobian)
print(f"UCM variance (per DOF): {v_ucm:.6f}")
print(f"ORT variance (per DOF): {v_ort:.6f}")
print(f"UCM/ORT ratio:          {ratio:.2f}")
print(f"(Ratio > 1 supports the UCM hypothesis)")
\end{lstlisting}

\section{Bernstein's Three Stages of Learning}

Bernstein proposed that motor skill acquisition progresses through three stages:

\begin{enumerate}
  \item \textbf{Freezing}: initially, the learner constrains (``freezes'') excess
        degrees of freedom, reducing the system to manageable dimensionality.
        A beginner golfer keeps wrists rigid, eliminating the wrist DOF.
  \item \textbf{Freeing}: as skill develops, the learner gradually releases
        (``frees'') the frozen DOF, allowing them to participate in the movement.
        The intermediate golfer adds wrist hinge.
  \item \textbf{Exploiting}: the expert learner exploits the released DOF to
        harness reactive forces, elastic energy, and passive dynamics. The expert
        golfer uses the passive wrist release driven by centrifugal force.
\end{enumerate}

In the language of control theory:
\begin{itemize}
  \item Stage 1: reduce the system dimension from $n$ to $k$ (the task dimension)
  \item Stage 2: increase the controlled dimension while maintaining task performance
  \item Stage 3: exploit the drift dynamics $f(\state)$ by designing trajectories that
        use passive mechanics (Volume~II, Chapter~5)
\end{itemize}

\section*{Exercises}

\begin{enumerate}
  \item \textbf{Redundancy Quantification.}
        For each of the following tasks, estimate the task dimension $k$, body DOF $n$,
        and compute the kinematic redundancy:
        (a) Standing balance, (b) Picking up a cup, (c) Walking on a treadmill,
        (d) A tennis serve.

  \item \textbf{UCM Analysis.}
        Using the provided code, generate synthetic data for 3-joint planar reaching
        with structured variability (more noise in the null space of the Jacobian).
        Verify that the UCM/ORT ratio exceeds 1.

  \item \textbf{Bernstein's Stages.}
        Observe a beginner and an expert performing the same motor task (e.g., throwing
        a ball). Describe their movement in terms of Bernstein's three stages. Which
        degrees of freedom are frozen, freed, and exploited?

  \item \textbf{(Programming.)} Implement a UCM analysis for a 7-DOF arm model reaching
        to a 3D target. Visualize the UCM and ORT directions.

  \item \textbf{(Research.)} Read Scholz and Sch{\"o}ner (1999) and summarize the key
        experimental evidence for the UCM hypothesis.
\end{enumerate}

\chapter{The Curse of Dimensionality --- And How Humans Solve It}
\label{ch:curse-dim}

\epigraph{%
  In ten dimensions, the unit hypercube has $2^{10} = 1024$ corners.
  In a hundred, it has $2^{100} \approx 10^{30}$.
  The nervous system controls hundreds of dimensions every second.
}{}

\section{Classical Curse of Dimensionality}

The curse of dimensionality \cite{Bellman1961} states that the computational resources
needed to solve many problems grow \emph{exponentially} with the state dimension $n$:

\begin{definition}{The Curse of Dimensionality}{curse-def}
  For a control problem with state dimension $n$:
  \begin{itemize}
    \item \textbf{Grid-based methods} (dynamic programming): cost $\sim \mathcal{O}(N^n)$
          where $N$ is the grid resolution per dimension
    \item \textbf{Sample-based methods} (Monte Carlo, RL): sample complexity grows
          polynomially in $n$ but with large constants
    \item \textbf{Trajectory optimization} (DDP, collocation): cost $\sim \mathcal{O}(n^3 \cdot T)$
          per iteration --- the only scalable approach for deterministic systems
  \end{itemize}
\end{definition}

\begin{figure}[h]
\centering
\begin{tikzpicture}[scale=1.0]
  % Axes
  \draw[->, thick] (0, 0) -- (6, 0) node[align=center, anchor=north] {Dimension ($n$)};
  \draw[->, thick] (0, 0) -- (0, 5) node[align=center, anchor=east, rotate=90] {Computational Cost};

  % Grid-based (exponential)
  \draw[red, thick, smooth] plot coordinates {
    (1, 0.3) (1.5, 0.6) (2, 1.2) (2.5, 2.4) (3, 4.8)
  } node[anchor=west] {\small Grid: $O(N^n)$};

  % DDP (cubic)
  \draw[blue, thick, smooth] plot coordinates {
    (1, 0.2) (2, 0.5) (3, 1.0) (4, 1.5) (5, 2.5) (6, 3.5)
  } node[anchor=west] {\small DDP: $O(n^3)$};

  % Synergy reduction (flat)
  \draw[green, thick] (0, 0.3) -- (6, 0.3) node[align=center, anchor=west] {\small Synergies: $O(k^3)$};

  % Shaded region (infeasible)
  \fill[red, opacity=0.1] (4, 5) -- (6, 5) -- (6, 0) -- (4, 0) -- cycle;
  \draw (5, 3) node[anchor=center, align=center] {\small Curse\\ \small region};

  \draw (3, 4.2) node[align=center, anchor=center, font=\small] {Exponential growth};
\end{tikzpicture}
\caption{The curse of dimensionality: computational cost grows exponentially with dimension
for grid-based methods, cubically for DDP, and remains constant when using synergy-based
dimensionality reduction (where $k \ll n$).}
\label{fig:ch2:curse-viz}
\end{figure}

\section{Biological Solutions}

Humans solve the dimensionality problem through five complementary mechanisms:

\subsection{1. Synergy-Based Compression}
As introduced in Chapter~4: $\bm{u}(t) = W \bm{c}(t)$ reduces $p$-dimensional muscle
space to $k$-dimensional synergy space, with $k/p \approx 4/400 = 1\%$.

\subsection{2. Hierarchical Decomposition}
Multi-level processing decomposes the global problem into independent subproblems.

\subsection{3. Passive Dynamics Exploitation}
The drift field $f(\state)$ does useful work. The controller only corrects deviations.
This reduces the effective control dimension (Volume~II, passive dynamics).

\subsection{4. Internal Models and Prediction}
Forward models predict consequences, enabling open-loop execution with sparse
feedback corrections (see Chapter~6).

\subsection{5. Task-Space Control}
Control in a low-dimensional task space rather than the full joint/muscle space.
The null space absorbs variability (UCM framework, Chapter~1).

\begin{lstlisting}[language=Python, caption={Computational cost comparison across methods.}]
import numpy as np

def computational_cost_table(n_dims: list[int]) -> None:
    """Print computational cost for different methods vs dimension."""
    print(f"{'n_dim':>6} {'Grid (N=10)':>15} {'DDP (T=100)':>15} "
          f"{'Synergy (k=5)':>15}")
    print("-" * 55)
    for n in n_dims:
        grid_cost = 10**n if n <= 10 else float('inf')
        ddp_cost = n**3 * 100
        synergy_cost = 5**3 * 100  # After reduction to k=5
        print(f"{n:>6d} {grid_cost:>15.0e} {ddp_cost:>15.0e} "
              f"{synergy_cost:>15.0e}")

computational_cost_table([2, 5, 10, 20, 50, 100, 200])
\end{lstlisting}

\section*{Exercises}

\begin{enumerate}
  \item \textbf{Scaling Analysis.}
        Using the computational cost table, determine the maximum state dimension for
        which grid-based dynamic programming is feasible (assuming $10^{12}$ FLOPS).

  \item \textbf{Synergy Efficiency.}
        If a motor task uses 100 muscles with 5 synergies, by what factor does the
        synergy decomposition reduce the dimensionality of the control problem?

  \item \textbf{(Programming.)} Implement DDP for a 2-DOF, 5-DOF, and 10-DOF chain.
        Measure computation time per iteration. Verify cubic scaling in state dimension.
\end{enumerate}

\chapter{Neural Architecture of Motor Control}
\label{ch:neural-arch}

\epigraph{%
  The motor system is not one controller.
  It is an orchestra of specialized modules,
  each playing its part in the symphony of movement.
}{}

\section{Cortical Motor Areas}

The cerebral cortex contains several distinct motor areas, each with specialized roles:

\subsection{Primary Motor Cortex (M1)}

M1 (Brodmann area 4) is the final cortical output stage for voluntary movement.
Layer V contains the giant pyramidal (Betz) cells whose axons form the corticospinal
tract---the most direct neural pathway from cortex to spinal motor neurons.

Key properties of M1:
\begin{itemize}
  \item \textbf{Somatotopic organization}: body parts are mapped to cortical regions
        (the motor homunculus), but with overlapping representations
  \item \textbf{Population coding}: individual neurons have broad tuning; movement
        direction is encoded by the \emph{population vector} \cite{Georgopoulos1986}
  \item \textbf{Preparatory activity}: M1 neurons are active 100--200~ms before
        movement onset, during motor preparation
\end{itemize}

\subsection{Premotor Cortex (PMC) and Supplementary Motor Area (SMA)}

\begin{itemize}
  \item \textbf{PMC}: movement planning based on external cues (visually guided reaching)
  \item \textbf{SMA}: movement planning based on internal cues (self-initiated sequences)
  \item Both areas project to M1 and directly to the spinal cord
\end{itemize}

\section{Subcortical Motor Structures}

\subsection{Cerebellum: The Forward Model}

The cerebellum contains more neurons than the rest of the brain combined ($\sim$70 billion
granule cells). Its primary role in motor control is as a \textbf{forward model}:

\begin{neurobox}{The Cerebellum as a Forward Model}{cerebellum-forward}
  The cerebellum:
  \begin{enumerate}
    \item Receives a copy of the motor command (efference copy) via the pontine nuclei
    \item Receives sensory feedback via the inferior olive and mossy fibers
    \item Computes the \emph{predicted} sensory consequences of the command
    \item Compares prediction with actual sensory feedback
    \item Generates an error signal that updates the forward model (climbing fiber
          pathway, long-term depression of parallel fiber--Purkinje cell synapses)
  \end{enumerate}
  In the language of control theory (Volume~I):
  \begin{itemize}
    \item The cerebellum implements $\hat{\state}_{k+1} = f(\hat{\state}_k, \control_k)$
          --- the state prediction step of the Kalman filter
    \item The climbing fiber error signal implements the innovation update
    \item The granule cell layer provides the basis functions for the nonlinear mapping
  \end{itemize}
\end{neurobox}

\subsection{Basal Ganglia: Action Selection}

The basal ganglia implement an action selection mechanism through competitive
inhibition. The direct pathway facilitates a desired action; the indirect pathway
suppresses competing actions. This is analogous to a model predictive controller
evaluating multiple candidate trajectories and selecting the best one.

\section{Spinal Cord: The Local Controller}

The spinal cord is not merely a relay station. It contains:
\begin{itemize}
  \item \textbf{Motor neuron pools}: lower motor neurons that directly innervate muscles
  \item \textbf{Interneuron networks}: local circuits that implement reflexes, pattern
        generation, and coordination
  \item \textbf{Central pattern generators} (CPGs): oscillatory circuits for rhythmic
        movements (Chapter~8)
\end{itemize}

\section*{Exercises}

\begin{enumerate}
  \item \textbf{Population Vector.}
        Given neurons with preferred directions every $45^\circ$, compute the population
        vector for a movement at $30^\circ$. Plot individual neuron contributions.

  \item \textbf{Cerebellar Forward Model.}
        Implement a simple forward model using a lookup table. Show how climbing
        fiber errors update the model over successive trials.

  \item \textbf{(Research.)} Read Georgopoulos et al.\ (1986) and summarize the
        evidence for population coding in M1.
\end{enumerate}

\chapter{Shallow-but-Wide: Why Brains Are Parallel, Not Deep}
\label{ch:shallow-wide}

\epigraph{%
  The brain is not a deep learning system.
  It is a massively parallel, remarkably shallow one.
  This is not a bug --- it is the architecture that makes
  real-time motor control possible.
}{}

\section{The Speed Constraint}

The most fundamental constraint on neural computation is \textbf{speed}. Every
biological computation has a time cost determined by the physics of synaptic
transmission:

\begin{neurobox}{Neural Timing Constraints}{timing}
  \begin{itemize}
    \item \textbf{Synaptic delay}: 1--5~ms per synapse (chemical synapses)
    \item \textbf{Axonal conduction}: 1--100~m/s (depends on myelination and diameter)
    \item \textbf{Action potential duration}: $\sim$1~ms
    \item \textbf{Refactory period}: $\sim$1--2~ms absolute, $\sim$3--5~ms relative
    \item \textbf{Total per-layer delay}: $\sim$5--10~ms including integration time
  \end{itemize}
\end{neurobox}

These delays are not technological limitations to be overcome with faster hardware.
They are \emph{fundamental physical constraints} of electrochemical signaling. No
amount of biological optimization can reduce synaptic delay below $\sim$1~ms.

\subsection{The Depth Budget}

Given the timing constraints, we can compute the maximum depth of any neural
computation that must complete within a given time window:

\begin{definition}{The Neural Depth Budget}{depth-budget}
  For a motor task requiring a response within time $T_{\text{response}}$, the
  maximum number of serial processing layers is:
  \begin{equation}
    L_{\max} = \frac{T_{\text{response}}}{\tau_{\text{layer}}},
    \label{eq:ch4:depth-budget}
  \end{equation}
  where $\tau_{\text{layer}} \approx 5$--$10$~ms is the per-layer processing time.
  For typical motor control scenarios:
  \begin{center}
  \renewcommand{\arraystretch}{1.3}
  \begin{tabular}{@{}lrr@{}}
  \toprule
  \textbf{Task} & $T_{\text{response}}$ & $L_{\max}$ \\
  \midrule
  Spinal reflex & 30~ms & 3--6 \\
  Postural correction & 100~ms & 10--20 \\
  Voluntary movement initiation & 150--250~ms & 15--50 \\
  Golf swing correction (mid-downswing) & 50~ms & 5--10 \\
  Saccadic eye movement & 200~ms & 20--40 \\
  \bottomrule
  \end{tabular}
  \end{center}
  Compare this with modern deep learning: GPT-4 has $\sim$120 transformer layers;
  ResNet-152 has 152 layers. These run on silicon at $\sim$ns per operation.
  Biology cannot afford this depth at $\sim$ms per operation.
\end{definition}

\section{The Motor Cortex: Six Layers, Not Sixty}

The cerebral cortex, including the primary motor cortex (M1), has a remarkably
consistent architecture across mammalian species \cite{Mountcastle1997}:

\begin{enumerate}
  \item \textbf{Layer I}: molecular layer (mostly dendrites and axons, few cell bodies)
  \item \textbf{Layer II/III}: external pyramidal layer (cortico-cortical connections)
  \item \textbf{Layer IV}: internal granular layer (thalamic input)
  \item \textbf{Layer V}: internal pyramidal layer (\emph{output layer} --- corticospinal tract)
  \item \textbf{Layer VI}: multiform layer (feedback to thalamus)
\end{enumerate}

The signal path from sensory input (Layer IV) through cortical processing to motor
output (Layer V) passes through \emph{at most 3--4 synapses within the cortex itself}.
Adding thalamic relay and spinal motor neuron connections, the total sensorimotor loop
is approximately 8--12 synapses deep.

\begin{insight}{The Architecture Is the Algorithm}{arch-is-algo}
  In deep learning, we can stack layers because silicon is fast. In biology, we cannot.
  Instead, the brain uses:
  \begin{enumerate}
    \item \textbf{Width}: $\sim$10 billion cortical neurons, each receiving $\sim$10,000
          synaptic inputs. The total computational capacity is in the \emph{connections},
          not the depth.
    \item \textbf{Recurrence}: horizontal connections within layers create recurrent
          dynamics that effectively increase computational depth through temporal
          iteration (but at the cost of time).
    \item \textbf{Specialization}: different brain regions process different aspects
          of the motor task in parallel, communicating through thick fiber bundles.
    \item \textbf{Hierarchy}: a shallow hierarchy of cortex $\to$ subcortical structures
          $\to$ spinal cord $\to$ muscles, with each level operating at different
          timescales and abstraction levels.
  \end{enumerate}
\end{insight}

\section{Computational Implications}

The shallow-but-wide architecture has profound implications for motor control theory:

\subsection{Complex Computations Through Simple Local Rules}

Each neuron performs a relatively simple computation: weighted sum of inputs, passed
through a nonlinear activation function. The power of the system comes from the
\emph{collective behavior} of billions of these simple units operating in parallel.

\begin{lstlisting}[language=Python, caption={Comparison: deep vs. wide network for motor control.}]
import numpy as np
import time

def deep_network_forward(
    x: np.ndarray,
    weights: list[np.ndarray],
    biases: list[np.ndarray],
    layer_delay_ms: float = 5.0
) -> tuple[np.ndarray, float]:
    """Forward pass through a deep network with biological timing.

    Parameters
    ----------
    x : np.ndarray
        Input vector.
    weights : list of np.ndarray
        Weight matrices for each layer.
    biases : list of np.ndarray
        Bias vectors for each layer.
    layer_delay_ms : float
        Simulated per-layer delay in milliseconds.

    Returns
    -------
    tuple
        (output, total_latency_ms)
    """
    h = x
    total_latency = 0.0
    for W, b in zip(weights, biases):
        h = np.maximum(0, W @ h + b)  # ReLU activation
        total_latency += layer_delay_ms
    return h, total_latency


def wide_network_forward(
    x: np.ndarray,
    W_hidden: np.ndarray,
    b_hidden: np.ndarray,
    W_out: np.ndarray,
    b_out: np.ndarray,
    layer_delay_ms: float = 5.0
) -> tuple[np.ndarray, float]:
    """Forward pass through a wide (2-layer) network.

    Parameters
    ----------
    x : np.ndarray
        Input vector.
    W_hidden : np.ndarray
        Hidden layer weights (wide).
    b_hidden : np.ndarray
        Hidden layer biases.
    W_out : np.ndarray
        Output layer weights.
    b_out : np.ndarray
        Output biases.
    layer_delay_ms : float
        Per-layer delay.

    Returns
    -------
    tuple
        (output, total_latency_ms)
    """
    h = np.maximum(0, W_hidden @ x + b_hidden)
    y = W_out @ h + b_out
    total_latency = 2 * layer_delay_ms  # Only 2 layers
    return y, total_latency


# Compare architectures for a motor control task
# Input: 20-dim state, Output: 10-dim muscle activations
np.random.seed(42)
n_input, n_output = 20, 10
n_params_budget = 50000  # Same total parameter budget

# Deep network: 10 layers of width 50
deep_widths = [n_input] + [50] * 9 + [n_output]
deep_W = [np.random.randn(deep_widths[i+1], deep_widths[i]) * 0.1
           for i in range(len(deep_widths)-1)]
deep_b = [np.zeros(deep_widths[i+1]) for i in range(len(deep_widths)-1)]

# Wide network: 2 layers, width 2000
wide_hidden = 2000
W_h = np.random.randn(wide_hidden, n_input) * 0.1
b_h = np.zeros(wide_hidden)
W_o = np.random.randn(n_output, wide_hidden) * 0.1
b_o = np.zeros(n_output)

x = np.random.randn(n_input)

y_deep, latency_deep = deep_network_forward(x, deep_W, deep_b)
y_wide, latency_wide = wide_network_forward(x, W_h, b_h, W_o, b_o)

print(f"Deep network (10 layers x 50 wide):")
print(f"  Latency: {latency_deep:.0f} ms")
print(f"  Parameters: {sum(W.size + b.size for W, b in zip(deep_W, deep_b)):,}")

print(f"\nWide network (2 layers x 2000 wide):")
print(f"  Latency: {latency_wide:.0f} ms")
print(f"  Parameters: {W_h.size + b_h.size + W_o.size + b_o.size:,}")

print(f"\nLatency ratio: {latency_deep/latency_wide:.1f}x faster with wide network")
\end{lstlisting}

\subsection{Hierarchical Decomposition}

The brain does not solve the full-dimensional control problem in one computation.
Instead, it decomposes the problem hierarchically:

\begin{enumerate}
  \item \textbf{High level} (prefrontal cortex, $\sim$200~ms): task goal, strategy
  \item \textbf{Mid level} (premotor/SMA, $\sim$100~ms): movement plan, sequencing
  \item \textbf{Low level} (M1 + spinal cord, $\sim$30~ms): muscle activation patterns
\end{enumerate}

Each level operates on a compressed representation of the level below:

\begin{proposition}
  In a hierarchical motor control architecture with $L$ levels, each reducing
  dimensionality by a factor $r$, the effective state dimension at level $l$ is:
  \begin{equation}
    n_l = n_0 / r^l,
    \label{eq:ch4:hierarchical-dim}
  \end{equation}
  where $n_0$ is the full state dimension. The computational cost per level
  scales as $\mathcal{O}(n_l^2)$ for a single-layer network, giving total
  cost $\sum_{l=0}^{L} n_l^2 = n_0^2 \sum_{l=0}^{L} r^{-2l}$, which converges
  for $r > 1$. The curse of dimensionality is broken by hierarchy.
\end{proposition}

\subsection{Synergies as Dimensionality Reduction}

Muscle synergies (d'Avella and Bizzi 2005 \cite{dAvellaBizzi2005}) provide the
most direct evidence for dimensionality reduction in motor control. The activation
pattern of $p$ muscles is approximated as a linear combination of $k \ll p$ synergies:
\begin{equation}
  \bm{u}(t) = \sum_{i=1}^{k} c_i(t) \bm{w}_i = W \bm{c}(t),
  \label{eq:ch4:synergy-decomposition}
\end{equation}
where $W \in \Reals^{p \times k}$ is the synergy matrix (fixed spatial patterns)
and $\bm{c}(t) \in \Reals^k$ are the time-varying synergy activations.

Experimental evidence consistently shows that $k = 4$--$6$ synergies can explain
$> 90\%$ of the variance in muscle activation patterns across a wide range of
tasks \cite{Tresch2006, Ting2007}.

\begin{lstlisting}[language=Python, caption={Muscle synergy extraction using Non-negative Matrix Factorization.}]
import numpy as np

def extract_synergies(
    emg_data: np.ndarray,
    n_synergies: int,
    max_iter: int = 500,
    tol: float = 1e-6
) -> tuple[np.ndarray, np.ndarray, float]:
    """Extract muscle synergies using NMF.

    Parameters
    ----------
    emg_data : np.ndarray, shape (n_timepoints, n_muscles)
        Processed EMG data (non-negative).
    n_synergies : int
        Number of synergies to extract.
    max_iter : int
        Maximum NMF iterations.
    tol : float
        Convergence tolerance.

    Returns
    -------
    tuple
        (synergy_weights, synergy_activations, vaf)
        W: (n_muscles, n_synergies)
        C: (n_timepoints, n_synergies)
        vaf: variance accounted for (0-1)
    """
    n_t, n_m = emg_data.shape

    # Initialize with random non-negative values
    W = np.abs(np.random.randn(n_m, n_synergies)) * 0.01
    C = np.abs(np.random.randn(n_t, n_synergies)) * 0.01

    for iteration in range(max_iter):
        # Multiplicative update rules (Lee & Seung 2001)
        # Update C
        numerator = emg_data @ W
        denominator = C @ W.T @ W + 1e-10
        C *= numerator / denominator

        # Update W
        numerator = emg_data.T @ C
        denominator = W @ C.T @ C + 1e-10
        W *= numerator / denominator

        # Check convergence
        reconstruction = C @ W.T
        residual = np.sum((emg_data - reconstruction)**2)
        total = np.sum(emg_data**2)
        vaf = 1.0 - residual / total

        if iteration > 0 and abs(vaf - prev_vaf) < tol:
            break
        prev_vaf = vaf

    return W, C, vaf


# Example: simulate EMG from 8 muscles during walking
np.random.seed(42)
n_timepoints = 200
n_muscles = 8

# Ground truth: 3 synergies generating 8-muscle activation
true_W = np.array([
    [0.8, 0.6, 0.1, 0.0, 0.2, 0.0, 0.1, 0.0],  # Syn 1: extensors
    [0.0, 0.1, 0.7, 0.9, 0.0, 0.3, 0.0, 0.1],  # Syn 2: flexors
    [0.2, 0.0, 0.1, 0.0, 0.8, 0.7, 0.6, 0.9],  # Syn 3: stabilizers
]).T  # (8 muscles, 3 synergies)

t = np.linspace(0, 2*np.pi, n_timepoints)
true_C = np.column_stack([
    np.maximum(0, np.sin(t)),
    np.maximum(0, np.sin(t + 2*np.pi/3)),
    np.maximum(0, np.sin(t + 4*np.pi/3)),
])

emg_simulated = true_C @ true_W.T + np.abs(np.random.randn(n_timepoints, n_muscles)) * 0.05

# Extract synergies
W_est, C_est, vaf = extract_synergies(emg_simulated, n_synergies=3)
print(f"Synergy extraction: {vaf:.1%} variance accounted for with 3 synergies")

# Test with different numbers
for k in range(1, 7):
    _, _, vaf_k = extract_synergies(emg_simulated, n_synergies=k)
    print(f"  k={k}: VAF = {vaf_k:.1%}")
\end{lstlisting}

\section{The Shallow-Wide Hypothesis for Motor Control}

We now formalize the central thesis of this volume:

\begin{insight}{The Shallow-Wide Hypothesis}{shallow-wide-hypothesis}
  Biological motor control systems achieve high-dimensional control through
  architectures that are:
  \begin{enumerate}
    \item \textbf{Shallow}: at most 10--20 serial processing layers (limited by
          synaptic delay)
    \item \textbf{Wide}: millions to billions of units per layer (enabled by neural
          density and parallel wiring)
    \item \textbf{Modular}: functionally specialized subnetworks process different
          aspects of the task in parallel
    \item \textbf{Hierarchical}: a small number of levels, each compressing dimensionality
    \item \textbf{Recurrent}: intra-layer connections provide temporal depth without
          adding spatial depth
  \end{enumerate}
  This architecture can approximate any smooth mapping (universal approximation theorem)
  with a single hidden layer of sufficient width. The practical advantage over deep
  networks is \textbf{latency}: a 2-layer network with 10,000 neurons per layer
  completes in 10~ms. A 100-layer network with 100 neurons per layer completes in
  500~ms---too slow for motor control.

  \textbf{For real-time motor control, width is cheap and depth is expensive.}
\end{insight}

\section*{Exercises}

\begin{enumerate}
  \item \textbf{Depth Budget.}
        For a goalkeeper diving to save a penalty kick (reaction time $\sim$200~ms),
        compute the maximum neural depth. How deep can the ``controller'' be?

  \item \textbf{Width vs. Depth.}
        Using the provided code, compare the outputs of a 10-layer $\times$ 50-wide
        network with a 2-layer $\times$ 2000-wide network on random inputs. Compare
        latencies.

  \item \textbf{Synergy Analysis.}
        Using the NMF code, determine how many synergies are needed to explain
        $> 95\%$ of variance in the simulated EMG data. Does this match the
        literature value of 4--6?

  \item \textbf{(Programming.)} Implement a motor control task (reaching) using both
        a deep and a wide network architecture. Compare tracking accuracy at different
        latency budgets.

  \item \textbf{(Research.)} Find experimental evidence for the number of cortical
        layers involved in processing a motor command from M1 to spinal motor neurons.
        How many synapses does the signal cross?
\end{enumerate}

\chapter{Ideomotor Theory and the Predictive Brain}
\label{ch:ideomotor}

\epigraph{%
  Every representation of a movement awakens in some degree
  the actual movement which is its object.
}{--- William James, 1890}

\section{Historical Ideomotor Theory}

The ideomotor principle, articulated by William James in 1890 \cite{James1890}, states
that the mere \emph{idea} of a movement tends to produce that movement. Modern
reformulations by Greenwald (1970) and Hommel et al.\ (2001) have transformed this
philosophical observation into a testable scientific framework.

\begin{definition}{The Theory of Event Coding (TEC)}{tec-def}
  The Theory of Event Coding \cite{Hommel2001} proposes that:
  \begin{enumerate}
    \item Perception and action share a \textbf{common representational format}
          (event codes or feature codes)
    \item Actions are planned in terms of their \textbf{desired sensory effects},
          not in terms of muscle commands
    \item The mapping from desired effects to motor commands is learned through
          experience (bidirectional action-effect associations)
  \end{enumerate}
\end{definition}

\section{The Predictive Processing Framework}

Modern neuroscience interprets the brain as a \textbf{prediction machine}
\cite{Clark2013, Hohwy2013}:

\begin{insight}{The Brain as a Prediction Machine}{prediction-machine}
  The brain continuously generates predictions about incoming sensory data.
  Neural processing is dominated by \textbf{prediction error minimization}:
  \begin{equation}
    \varepsilon_k = \bm{y}_k - \hat{\bm{y}}_k = \bm{y}_k - h(\hat{\state}_k),
    \label{eq:ch5:prediction-error}
  \end{equation}
  where $\hat{\bm{y}}_k$ is the predicted observation and $\bm{y}_k$ is the actual
  observation. This error drives updates to the internal model.

  In the language of Volume~I, Chapter~6: this is exactly the
  \textbf{Kalman filter innovation}. The brain implements Bayesian state estimation.
\end{insight}

\subsection{Active Inference}

Karl Friston's active inference framework \cite{Friston2010} extends predictive
processing to action: instead of only updating beliefs to match observations,
the agent can also \emph{act} to make observations match predictions:
\begin{equation}
  \control^* = \argmin_{\control} \mathbb{E}\left[
  \sum_{k=0}^{T} \norm{\bm{y}_k - \hat{\bm{y}}_k}^2_{R^{-1}}
  + \norm{\control_k}^2_Q
  \right],
  \label{eq:ch5:active-inference}
\end{equation}
which is formally equivalent to model predictive control (Volume~I, Chapter~5).

\section{Connection to Control Theory}

\begin{center}
\renewcommand{\arraystretch}{1.3}
\begin{tabular}{@{}ll@{}}
\toprule
\textbf{Neuroscience Term} & \textbf{Control Theory Term} \\
\midrule
Forward model & State transition matrix / predictor \\
Prediction error & Innovation (Kalman filter) \\
Efference copy & Feedforward control signal \\
Active inference & Model predictive control \\
Prior beliefs & State constraints / initial conditions \\
Precision weighting & Kalman gain / observation weights \\
Free energy & Variational cost functional \\
\bottomrule
\end{tabular}
\end{center}

\section*{Exercises}

\begin{enumerate}
  \item \textbf{Conceptual.} Explain how the ideomotor principle relates to
        trajectory-centric control (Volume~II). When a golfer imagines impact,
        what ``prediction'' are they generating?

  \item \textbf{Kalman Connection.} Write out the Kalman filter equations and identify
        which term corresponds to the prediction error in the predictive processing
        framework.

  \item \textbf{(Programming.)} Implement an active inference agent that controls a
        pendulum by minimizing prediction error rather than tracking a reference.
\end{enumerate}

\chapter{Internal Models: Forward and Inverse}
\label{ch:internal-models}

\epigraph{%
  To control is to predict.
  To predict is to have a model.
  To have a model is to have learned.
}{}

\section{Forward Models}

A forward model predicts the sensory consequences of a motor command:
\begin{equation}
  \hat{\state}_{k+1} = f(\hat{\state}_k, \control_k),
  \label{eq:ch6:forward-model}
\end{equation}
which is exactly the state propagation used in Volume~I.

\subsection{The MOSAIC Architecture}

The MOSAIC model (Modular Selection and Identification for Control)
\cite{WolpertKawato1998} proposes that the brain maintains \emph{multiple}
paired forward-inverse models, each tuned to a different context:

\begin{lstlisting}[language=Python, caption={MOSAIC-style multiple model architecture.}]
import numpy as np
from typing import List, Tuple

class ForwardInversePair:
    """A paired forward-inverse model for one context."""
    def __init__(self, A: np.ndarray, B: np.ndarray):
        self.A = A  # State transition
        self.B = B  # Input matrix

    def forward_predict(self, x: np.ndarray,
                        u: np.ndarray) -> np.ndarray:
        """Predict next state."""
        return self.A @ x + self.B @ u

    def inverse(self, x: np.ndarray,
                x_desired: np.ndarray) -> np.ndarray:
        """Compute control to reach desired state."""
        return np.linalg.lstsq(self.B,
                               x_desired - self.A @ x,
                               rcond=None)[0]


class MOSAIC:
    """Multiple paired forward-inverse model architecture."""
    def __init__(self, models: List[ForwardInversePair]):
        self.models = models
        self.n_models = len(models)
        self.responsibilities = np.ones(self.n_models) / self.n_models

    def update_responsibilities(self, x: np.ndarray,
                                u: np.ndarray,
                                x_next: np.ndarray) -> None:
        """Update model responsibilities based on prediction errors."""
        errors = np.zeros(self.n_models)
        for i, model in enumerate(self.models):
            prediction = model.forward_predict(x, u)
            errors[i] = np.sum((x_next - prediction)**2)

        # Softmax weighting (lower error = higher responsibility)
        log_resp = -0.5 * errors
        log_resp -= np.max(log_resp)  # numerical stability
        self.responsibilities = np.exp(log_resp)
        self.responsibilities /= np.sum(self.responsibilities)

    def predict(self, x: np.ndarray,
                u: np.ndarray) -> np.ndarray:
        """Weighted prediction across all models."""
        prediction = np.zeros_like(x)
        for i, model in enumerate(self.models):
            prediction += self.responsibilities[i] * \
                         model.forward_predict(x, u)
        return prediction

    def control(self, x: np.ndarray,
                x_desired: np.ndarray) -> np.ndarray:
        """Weighted inverse model output."""
        u = np.zeros(self.models[0].B.shape[1])
        for i, model in enumerate(self.models):
            u += self.responsibilities[i] * \
                 model.inverse(x, x_desired)
        return u
\end{lstlisting}

\section{Inverse Models}

An inverse model computes the motor command required to achieve a desired state:
\begin{equation}
  \control_k = f^{-1}(\state_k, \state_{k+1}^{\text{des}}),
  \label{eq:ch6:inverse-model}
\end{equation}
which corresponds to the feedforward control of Volume~II.

\section{Smith Predictor for Neural Delays}

The sensorimotor system faces delays of 50--200~ms. The Smith predictor
architecture uses a forward model to compensate:
\begin{equation}
  \hat{\state}_{\text{current}} = \hat{\state}(t - \Delta t) +
  \int_{t-\Delta t}^{t} f(\hat{\state}, \control) \, \dd\tau,
  \label{eq:ch6:smith-predictor}
\end{equation}
where $\Delta t$ is the sensory delay. The controller operates on the
\emph{predicted current state} rather than the delayed sensory measurement.

\section*{Exercises}

\begin{enumerate}
  \item \textbf{MOSAIC.} Using the provided code, simulate a system that switches
        between two dynamics (e.g., carrying an empty cup vs.\ a full cup). Show
        how the responsibilities shift.

  \item \textbf{Smith Predictor.} Implement a Smith predictor for a simple pendulum
        with a 100~ms sensory delay. Compare performance with and without prediction.

  \item \textbf{(Research.)} Summarize evidence for multiple internal models in
        human motor control, citing at least three experimental studies.
\end{enumerate}

\chapter{Passive Distributed Control}
\label{ch:passive-control}

\epigraph{%
  The most efficient movement is one where the body's own mechanics
  do most of the work, and the nervous system merely guides.
}{--- Adapted from Frans Bosch}

\section{The Frans Bosch Framework}

Frans Bosch's work on strength training and coordination \cite{Bosch2015}
provides a practitioner-oriented framework for understanding how elite athletes
organize movement. His key insights, translated into the language of this series:

\begin{insight}{Self-Organization in Complex Movement}{bosch-self-org}
  Bosch argues that complex movements are not centrally programmed but
  \textbf{self-organize} through the interaction of:
  \begin{enumerate}
    \item \textbf{Task constraints}: the goal shapes the solution space
    \item \textbf{Organismic constraints}: the body's anatomy and current state
    \item \textbf{Environmental constraints}: gravity, ground, implements
  \end{enumerate}
  The role of the nervous system is not to specify every muscle activation but to
  set the \emph{initial conditions} and \emph{constraints} that allow self-organization
  to produce the desired movement.

  In the language of Volume~I: the CNS designs the \emph{drift field} $f(\state)$
  of the system (through postural configuration and muscle tone) and lets the
  natural dynamics do the work.
\end{insight}

\section{Passive Dynamic Walking}

The most striking example of passive dynamics is passive dynamic walking
(McGeer, 1990): a simple mechanical biped with knees, feet, and hip (no motors)
walks down a gentle slope in stable limit cycles with zero active control. This
demonstrates that complex coordinated movement can emerge from mechanical design
and gravity alone.

\noindent
Mathematically, the walker is a system with zero control input: $\dot{\state} = f(\state)$
is purely drift dynamics with no control term $G(\state) \bm{u}$. The stability analysis
of periodic orbits in such systems is a central topic in Agrachev \& Sachkov \cite{AgrachevSachkov2004},
Chapter~2 (affine systems with zero control) and Chapter~4 (limit cycles and stability):
the existence of stable walking gaits relies on eigenvalue analysis of the Poincaré return map,
a classical dynamical systems tool that forms the foundation of their geometric approach.
The most striking example of passive dynamics is passive dynamic walking,
demonstrated by McGeer (1990) \cite{McGeer1990}: a simple mechanical biped
can walk down a gentle slope with \emph{no actuators and no control}:

\begin{lstlisting}[language=Python, caption={Passive dynamic walker simulation.}]
import numpy as np
from scipy.integrate import solve_ivp

def passive_walker_dynamics(
    t: float,
    state: np.ndarray,
    M: float = 10.0,
    m: float = 5.0,
    l: float = 1.0,
    g: float = 9.81,
    slope: float = 0.05
) -> np.ndarray:
    """Dynamics of a 2-DOF passive walker (compass gait).

    Parameters
    ----------
    t : float
        Time.
    state : np.ndarray
        [theta_stance, theta_swing, d_theta_stance, d_theta_swing]
    M : float
        Hip mass (kg).
    m : float
        Leg mass at foot (kg).
    g : float
        Gravitational acceleration (m/s^2).
    slope : float
        Ground slope (radians).

    Returns
    -------
    np.ndarray
        State derivatives.
    """
    th_st, th_sw, dth_st, dth_sw = state

    # Mass matrix
    M11 = (M + m) * l**2
    M12 = -m * l**2 * np.cos(th_st - th_sw)
    M21 = M12
    M22 = m * l**2

    # Gravity and Coriolis
    C1 = -m * l**2 * dth_sw**2 * np.sin(th_st - th_sw) \
         - (M + m) * g * l * np.sin(th_st - slope)
    C2 = m * l**2 * dth_st**2 * np.sin(th_st - th_sw) \
         + m * g * l * np.sin(th_sw - slope)

    # Solve M * ddq = C (NO control input!)
    det = M11 * M22 - M12 * M21
    ddth_st = (M22 * C1 - M12 * C2) / det
    ddth_sw = (M11 * C2 - M21 * C1) / det

    return np.array([dth_st, dth_sw, ddth_st, ddth_sw])


# Simulate one step of passive walking
state0 = np.array([0.2, -0.4, -0.5, 1.0])
sol = solve_ivp(passive_walker_dynamics, [0, 1.5], state0,
                max_step=0.001, events=None)
print(f"Passive walker: simulated {sol.t[-1]:.2f}s "
      f"with {len(sol.t)} timesteps")
print(f"Final stance angle: {sol.y[0, -1]:.3f} rad")
print("(No actuators, no controller --- just gravity and mechanics)")
\end{lstlisting}

\section{Passive Dynamics in the Golf Swing}

The golf swing exemplifies passive dynamics exploitation:
\begin{enumerate}
  \item \textbf{Backswing}: energy stored in elastic tissues (muscles, tendons, fascia)
  \item \textbf{Transition}: elastic recoil initiates the downswing passively
  \item \textbf{Wrist release}: centrifugal force passively uncocks the wrists
        (the double pendulum effect from Volume~II, Chapter~5)
  \item \textbf{Follow-through}: passive deceleration through tissue compliance
\end{enumerate}

In the ZTCF framework of Volume~I, Chapter~7: the drift contribution $f(\state)$
accounts for $\sim$60\% of the clubhead speed. The control contribution
$G(\state)\control$ provides the remaining $\sim$40\%.

Trajectory design exploiting drift is formalized in Agrachev \& Sachkov \cite{AgrachevSachkov2004},
Chapter~5 (optimal control with drift): when the drift term dominates the control term, optimal
trajectories naturally evolve along low-cost paths in the accessible set, allowing the system to
reach distant targets with minimal control effort. The golf swing is a biological instance of
this principle.

\section{Bernstein's Stage 3: Exploiting DOF}

The connection to Bernstein's learning framework (Chapter~1) is direct:
\begin{itemize}
  \item The expert athlete has reached Stage 3 --- exploiting degrees of freedom
  \item They have learned to configure the body so that passive dynamics
        \emph{automatically} produce the desired motion
  \item The control problem is simplified: instead of commanding every muscle,
        set up the initial conditions and let physics do the work
\end{itemize}

\section*{Exercises}

\begin{enumerate}
  \item \textbf{Passive Walking.} Using the provided code, find the set of
        initial conditions that produce a stable limit cycle (periodic gait).

  \item \textbf{Energy Analysis.} For the passive walker, plot kinetic and
        potential energy over one stride. Where does the energy come from?

  \item \textbf{Wrist Release.} Model the wrist release as a passive double
        pendulum. For what initial angular velocity of the arm does the club
        reach maximum speed at the bottom?

  \item \textbf{(Programming.)} Add a small actuator to the passive walker
        and design a minimal controller that stabilizes the gait on level ground.
\end{enumerate}

\chapter[Biology and Nonlinear Dynamics]{The Interplay of Biology and\\Dynamics of Nonlinear Systems}
\label{ch:biology-nonlinear}

\section{Introduction: Mechanics With a Biological State}

A golfer can arrive at nearly the same joint configuration and velocity with
different muscle activation, tendon strain and sensory estimates. Those states
need not produce the same subsequent motion. Geometry and rigid-body mechanics
remain essential, but positions and velocities alone may no longer close the
model. The additional state determines how a neural command becomes force, how
a disturbance propagates and which responses remain feasible before impact.

This is the useful connection between biology and nonlinear dynamics. Tissue
properties, neural feedback, task goals and contact conditions interact through
specific equations. Their interaction can support coordination; it can also
produce delay, instability, fatigue and estimation error. A biological label
does not decide which outcome occurs.

\begin{insight}[colframe=chapblue]{Three Questions That Must Be Kept Together}{bio-levels}
\begin{enumerate}
\item \textbf{Mechanics}: what forces and moments follow from the current
geometry, motion, tissue state and contact conditions?
\item \textbf{Control}: which excitation or actuator histories can change that
state within the remaining time and input limits?
\item \textbf{Evidence}: which measurements distinguish the proposed mechanism
from another model that produces the same observed movement?
\end{enumerate}
For a swing, the endpoint objective includes the clubhead's delivered position,
orientation and velocity, followed by collision dynamics. For a humanoid holding
a tool, it also includes grasp feasibility and balance. Neither objective is
specified by a single joint torque or a single measure of movement variability.
\end{insight}

\subsection{Bernstein, Constraints and the Scope of the Claim}

Having more muscle forces than independent joint moments makes force allocation
underdetermined unless additional information or assumptions are supplied. It
does not prove that the nervous system cannot coordinate the muscles, nor that
it must use a fixed small number of commands. A model's degrees of freedom,
actuators, measured EMG channels and task variables are different counts.

Newell grouped constraints into organism, task and environment~\cite{Newell1986}.
Bosch discusses stable and adaptable movement components in
sport~\cite{Bosch2020,Bosch2015}.
These frameworks motivate hypotheses; formal stability results require a
separate model and proof. Bosch's ``fluctuations''
can denote flexible task-dependent components; they are not, by definition,
expensive errors that must always be suppressed.

An intersection of feasible sets constrains possible states. It does not by
itself specify a vector field, an attractor or a controller. Two systems with
identical feasible configurations can have different inertia, dissipation and
feedback, and therefore different trajectories. Constraints-led practice and
predictive or feedback control are compatible descriptions at different levels.
Whether a golfer uses a particular internal representation requires evidence
beyond a successful dynamical model of the movement.

\section{Muscle, Tendon and Activation Dynamics}

\subsection{A Consistent Hill-Type Model}

Hill's shortening relation differs from later Hill-type muscle--tendon
models~\cite{Hill1938,millard2013flexing}. Declare which elements
are included, where their lengths are measured and how force is transmitted.
In an illustrative massless equilibrium model, a contractile element and passive
parallel element share the fiber length $l_f$. Their force is transmitted through
a series tendon, with pennation angle $\alpha$:
\begin{align}
F_f &= a F_0 f_L(l_f/l_0) f_V(\dot l_f/v_{\max})+F_{\mathrm{PE}}(l_f),
\label{eq:ch8:force-production}\\
F_t &= F_f\cos\alpha, \qquad
l_{\mathrm{MT}}(q)=l_t+l_f\cos\alpha.
\label{eq:ch8:hill-model}
\end{align}
Here $a\in[0,1]$ is activation, $F_0$ the reference active isometric fiber force,
$l_0$ the optimal fiber length, and $v_{\max}>0$ a shortening-speed scale.
The parallel force uses fiber length, not the total muscle--tendon length.
Pennation requires its own geometric law if it changes. Fiber and tendon forces
are not added as though series elements were parallel force sources.

A simple tension-only tendon law is
\begin{equation}
F_t=k_t[l_t-l_{\mathrm{slack}}]_+,
\qquad [z]_+=\max(z,0),
\label{eq:ch8:tendon-law}
\end{equation}
with $k_t$ in N/m. The extension is relative to tendon slack length, not muscle
length. This piecewise linear law is a teaching approximation; physiological
tendon models can include a nonlinear toe region and hysteresis. Series
equilibrium and path length together determine fiber velocity or require an
implicit solve. Slack tendon and near-zero activation need careful numerical
treatment rather than division by a vanishing active-force term.

\begin{figure}[htbp]
\centering
\begin{tikzpicture}[x=0.85cm,y=0.85cm]
\draw (0,0)--(1,0)--(1,0.8)--(2,0.8);
\draw (2,0.45) rectangle (3.2,1.15);
\node at (2.6,0.8) {CE};
\draw (3.2,0.8)--(4,0.8)--(4,0);
\draw (1,0)--(1,-0.8)--(2,-0.8);
\draw (2,-1.15) rectangle (3.2,-0.45);
\node at (2.6,-0.8) {PE};
\draw (3.2,-0.8)--(4,-0.8)--(4,0)--(4.6,0);
\draw (4.6,-0.35) rectangle (6.2,0.35);
\node at (5.4,0) {Tendon};
\draw (6.2,0)--(7,0);
\draw[<->] (1,-1.6)--(4,-1.6);
\node[below] at (2.5,-1.6) {$l_f$};
\draw[<->] (4,-1.6)--(7,-1.6);
\node[below] at (5.5,-1.6) {$l_t$};
\draw[->,thick] (7,0)--(7.7,0);
\node[above] at (7.3,0) {$F_t$};
\end{tikzpicture}
\caption{Parallel Fiber Elements and a Series Tendon. The Schematic Uses Zero
Pennation; a Pennate Model Must Project Fiber Force and Length Consistently.}
\label{fig:ch8:hill-model}
\end{figure}

\subsection{Length and Velocity Conventions That Can Be Checked}

Let $\nu=\dot l_f/v_{\max}$, so shortening means $\nu<0$.
One explicit illustrative pair of normalized curves is
\begin{align}
f_L(\ell)&=\exp\left[-\left(\frac{\ell-1}{w}\right)^2\right],\quad w>0,
\label{eq:ch8:force-length}\\
f_V(\nu)&=
\begin{cases}
\dfrac{1+\nu}{1-\nu/c},&-1\leq\nu\leq0,\\[4pt]
1+(N-1)\dfrac{\nu}{d+\nu},&\nu>0,
\end{cases}
\quad c,d>0,\quad N>1.
\label{eq:ch8:force-velocity}
\end{align}
These equations define the example, not a universal parameter fit or Hill's
original eccentric law. They give $f_V(-1)=0$, $f_V(0)=1$, and an eccentric
limit $N$. The shortening branch increases with $\nu$; faster shortening
therefore reduces force. Its derivative at zero is $1+1/c$. Choosing
$d=(N-1)/(1+1/c)$ also matches the eccentric branch's derivative there.
Outside the stated domain, choose and validate an extension rather than letting
negative active tension silently enter a simulation.

For $c=0.25$, $N=1.5$ and $d=0.1$, $f_V(-0.5)=1/6$ and
$f_V(0.1)=1.25$. At $a=0.5$, $F_0=1000$ N and optimal length, active fiber
force is respectively $83.33$ N and $625$ N. These are chosen model numbers,
not measured human values. Passive force and pennation still affect tendon
force. Total force at a long fiber length can exceed active isometric force
because the passive element contributes; the active force--length peak alone
does not bound the total tension.

\subsection{Zero Excitation Retains the Activation History}

Excitation $u$ and activation $a$ are different variables. For a constant
activation time scale $\tau_a>0$, a simple example is
\begin{equation}
\dot a=(u-a)/\tau_a,\qquad 0\leq u\leq1.
\label{eq:ch8:activation}
\end{equation}
If $u$ is set to zero at $t_0$, then
$a(t)=a(t_0)\exp[-(t-t_0)/\tau_a]$. Active force generally persists while
activation decays. Real activation/deactivation dynamics can require different
time scales and nonlinearities. With activation, fiber length and other internal
states $z$, a closed model uses $x=(q,v,a,z)$, not only $x=(q,v)$.

For muscle path lengths $l_{\mathrm{MT}}(q)$, define
$J_l=\partial l_{\mathrm{MT}}/\partial q$. Positive tendon tension pulls toward
shorter path length, giving generalized muscle load $-J_l^T F_t$:
\begin{align}
M(q)\dot v+h(q,v)&=-J_l(q)^TF_t(x)+Q_{\mathrm{ext}}+A(q)^T\lambda,
\label{eq:ch8:body-dynamics}\\
\dot q&=v,\qquad \dot z=\psi(x),\qquad
Av=0.
\end{align}
Here generalized coordinates are chosen so $\dot q=v$; other representations
need their kinematic map. $h$ includes the declared gravity and velocity terms,
and $\lambda$ enforces the current independent ideal constraints. Ground friction,
grasp limits and changing contact modes need additional conditions.

With constant $\tau_a$ and no other direct excitation dependence, the augmented
system is control-affine: $\dot x=f(x)+G(x)u$. Excitation enters the activation
rows; its influence on acceleration develops through the state. The drift
$f$ at $u=0$ includes existing activation, tendon strain and constraint response.
Holding a nonzero baseline command instead defines another vector field. It is
legitimate to study that field, but it must not be called fully passive or free
of ongoing energetic cost.

\subsection{Preflex Response and Tendon Energy}

An intrinsic mechanical response can precede a newly evoked neural response.
At a fixed activation and around a specified state, a local approximation is
\begin{equation}
\delta F_f\simeq K_f\delta l_f+D_f\delta\dot l_f.
\label{eq:ch8:preflex-force}
\end{equation}
The coefficients depend on operating point and history, not just activation.
This approximation is often associated with the preflex concept~\cite{Loeb1995}.
It introduces no new neural-loop delay; that does not imply instantaneous global
force transmission, an inactive nervous system, or guaranteed stability.
The sign and size of incremental stiffness, tendon compliance, limb geometry
and contact determine the resulting motion.

For example, a damped inverted pendulum near upright obeys
\begin{equation}
I\ddot\theta+D\dot\theta+(K-mgh)\theta=0.
\label{eq:ch8:posture-stability}
\end{equation}
With $I>0$ and $D>0$, this linear model is asymptotically stable only when
$K>mgh$. Positive tissue stiffness is insufficient if it is smaller than the
destabilizing gravitational term. Direct human standing measurements by Loram
and Lakie found intrinsic ankle stiffness insufficient for stability in their
protocol~\cite{LoramLakie2002}. This is a useful counterexample to a general
claim that baseline muscle tone automatically stabilizes posture. Spinal
circuits are part of the CNS; anesthesia does not imply elimination of every
reflex. Neither claim is needed to explain a mechanical preflex.

For the elastic law above,
\begin{equation}
E_t=\tfrac12 k_t[l_t-l_{\mathrm{slack}}]_+^2,
\qquad \dot E_t=F_t\dot l_t
\label{eq:ch8:tendon-energy}
\end{equation}
away from the switching point. A chosen $k_t=20{,}000$ N/m and 0.02 m
extension store 4 J at 400 N. Releasing 4 J in 0.02 s supplies 200 W of average
recoil power in this ideal example. The energy had to be loaded earlier;
recoil changes the timing of power without creating energy. Fiber work,
skeletal work, tendon storage and losses must close one balance. There is no
universal 50\% tendon contribution to jumping or striking, and an elastic
element is not generally equivalent to rescaling time in a rigid-body solution.

For a golf swing, a catapult hypothesis therefore needs a measured or bounded
strain history and release timing. A compliant robot actuator needs the same
accounting. More compliance can improve a specified objective or impair force
bandwidth, precision and contact control; the answer depends on the task.

\section{Attractors and Coordination Transitions}

\subsection{What Stability Does and Does Not Establish}

For a declared autonomous closed-loop system, we use an asymptotically stable
attractor to mean a compact invariant set that is Lyapunov stable and attracts
some neighborhood. Lyapunov stability means sufficiently small disturbances
remain small; attraction additionally concerns convergence to the set.
Definitions of more general attractors differ, so minimality and a single
negative exponent should not be substituted for these conditions.

Attraction to a periodic orbit is distinct from convergence to one time-stamped
trajectory on that orbit. A stable autonomous limit cycle has a neutral phase
direction, with a zero Lyapunov exponent and a unit Floquet multiplier; its
transverse multipliers can lie inside the unit circle. Nearby motions can retain
a phase offset while approaching the same orbit. A chaotic attractor can attract
trajectories to its set while separating nearby trajectories within that set.
Finite-time growth, asymptotic exponents and basin size answer different questions.

These distinctions matter in golf. A repeatable swing is a finite-duration
movement ending in contact, not automatically a limit cycle or equilibrium.
An appropriate analysis can use a trajectory tube, terminal error, input limits
and impact-event sensitivity. A phase error that is harmless for convergence to
a periodic gait can be consequential when club and ball must meet at one time.

An attractor can be generated by passive mechanics, constant powered actuation,
neural feedback or a robot controller. Its existence does not identify which
mechanism supplies stability, how much power is consumed or whether a new task
can be accommodated. A passive downhill walker draws energy from gravity;
an actively stabilized posture can consume energy without moving. A finite
feedforward maneuver need not be an attractor to meet an error tolerance.

\begin{example}[colframe=chapblue]{A Stable Orbit With a Neutral Phase}{bio-limit-cycle}
On an annulus around $r=1$, consider $\dot r=1-r$, $\dot\theta=\omega>0$.
Radial deviations decay as $e^{-t}$, but two solutions' angular difference
remains constant. The orbit $r=1$ attracts nearby states although their
time-matched Euclidean separation need not tend to zero. Transverse recovery
and phase recovery must therefore be measured separately.
\end{example}

\subsection{A Cost Crossing Is Not a Dynamical Bifurcation}

A dynamical bifurcation changes the qualitative behavior of a vector field or
return map as a parameter varies. Two metabolic cost curves crossing does not
show that either gait loses dynamical stability. Even in the simple chosen
costs $C_1(v)=1+v^2$ and $C_2(v)=2+v^2/2$, the crossing at $v=\sqrt2$
says only which algebraic cost is smaller. It supplies no equations for a gait.
An observed walk--run transition can be investigated using energetic, mechanical,
perceptual and stability hypotheses; matching one transition speed does not
distinguish them. These teaching curves are dimensionless examples, not fitted
human metabolic laws.

Likewise, a coordination variable and an experimentally adjusted ``control
parameter'' are not necessarily a plant state and actuator input in the
control-theory sense. The relationship must be derived or fitted. Frequency
may alter several physiological and mechanical processes at once.

\subsection{The Haken--Kelso--Bunz Model With Its Correct Transition}
\label{sec:bio-hkb}

In the classic bimanual experiment, increasing cycling frequency can cause the
initial anti-phase mode to switch to in-phase coordination, using the authors'
homologous-muscle phase convention~\cite{HakenKelsoBunz1985}. Symmetric movements
toward and away from the body's midline can be in-phase even though the hands'
world-coordinate velocities point in opposite directions. Define the measured
phase before using ``same direction'' as a substitute.

The standard reduced deterministic phase equation is first order:
\begin{align}
\dot\phi&=-a_H\sin\phi-2b_H\sin(2\phi)=-V'(\phi),
\label{eq:ch8:hkb-model}\\
V(\phi)&=-a_H\cos\phi-b_H\cos(2\phi),\qquad a_H,b_H>0.
\end{align}
The potential is a reduced dynamical construction, not metabolic or mechanical
energy. For fixed parameters, $\dot V=-[V'(\phi)]^2\leq0$. Linearizing at the
two coordination modes gives
\begin{equation}
\lambda_0=-a_H-4b_H<0,\qquad
\lambda_\pi=a_H-4b_H.
\label{eq:ch8:hkb-stability}
\end{equation}
Thus in-phase $\phi=0$ remains locally stable, whereas anti-phase $\phi=\pi$
is stable only for $b_H/a_H>1/4$. If the ratio decreases through $1/4$ as
cycling frequency increases, anti-phase loses stability. A chosen frequency
law is required before assigning a transition frequency in Hz. A forced
second-order phase oscillator is a different model requiring its own derivation.

For $a_H=1$ s$^{-1}$, changing $b_H$ from 0.5 to 0.2 s$^{-1}$ changes
$\lambda_\pi$ from $-1$ to $+0.2$ s$^{-1}$ while $\lambda_0$ remains negative.
Exactly $\phi=\pi$ is an equilibrium even after losing stability: a numerical
experiment must include a small perturbation or a declared noise process to
observe departure. Near the transition, slow recovery and noise-induced switching
can complicate identification of the deterministic threshold.

\begin{figure}[htbp]
\centering
\begin{tikzpicture}
\begin{axis}[width=0.95\linewidth,height=5.2cm,
xmin=0,xmax=0.5,ymin=-3.2,ymax=1.2,
xlabel={$b_H/a_H$},ylabel={$\lambda/a_H$},
xtick={0,0.125,0.25,0.375,0.5},ytick={-3,-2,-1,0,1},
legend style={at={(0.5,-0.3)},anchor=north,legend columns=2},
grid=major]
\addplot[blue,thick,domain=0:0.5] {-1-4*x};
\addlegendentry{In-Phase}
\addplot[red,thick,domain=0:0.5] {1-4*x};
\addlegendentry{Anti-Phase}
\addplot[black,dashed,forget plot] coordinates {(0,0) (0.5,0)};
\addplot[black,dotted,forget plot] coordinates {(0.25,-3.2) (0.25,1.2)};
\end{axis}
\end{tikzpicture}
\caption{Local Phase Stability in the HKB Reduction. Negative Rates Indicate
Recovery; Anti-Phase Loses Stability as the Ratio Falls Below One Quarter.}
\label{fig:ch8:hkb-bifurcation}
\end{figure}

The reduced model explains an experimental coordination transition without an
explicit discrete mode-selection command in its equations. It does not show
that neurons are absent, that the brain never predicts trajectories or that
golf timing follows the same phase law. In a swing, a proposed collective
variable needs a measurement definition, fitted dynamics and tests under
perturbations and altered tempo. A recognizable pattern is a starting point
for that work, not its causal conclusion.

\section{Synergies, Task Geometry and Control Authority}

\subsection{What a Low-Dimensional Fit Establishes}

A spatial muscle-synergy model approximates an EMG or activation matrix
$U\in\mathbb R^{p\times N}$ by $WC$, with $W\in\mathbb R^{p\times k}$ and
$C\in\mathbb R^{k\times N}$. Nonnegative factorization usually imposes
$W,C\geq0$. Specify preprocessing, normalization, muscles, tasks and rank
selection; EMG is not a direct measurement of muscle force.
\begin{equation}
\mathrm{VAF}_{\mathrm{uncentered}}
=100\left(1-\frac{\|U-WC\|_F^2}{\|U\|_F^2}\right)\%,
\qquad \|U\|_F>0.
\label{eq:ch8:vaf}
\end{equation}
The percentage multiplies the entire bracket. A centered variance definition
uses a different denominator and must be named. High reconstruction quality
is a property of the dataset and fit, not proof of a neural module or a
universal number of synergies.

Kutch and Valero-Cuevas supplied counterexamples in which biomechanics
produces low-dimensional activation patterns without fixed neural grouping
\cite{KutchValeroCuevas2012}. Such alternatives do not establish that neural
synergies never exist. They show why a causal claim needs task variation,
perturbation, held-out prediction and comparison with independently controlled
muscles. A useful robot policy may intentionally restrict its inputs to
synergies even when the corresponding human mechanism remains unresolved.

\subsection{A Task Level Set and Its Tangent Space}

For a smooth task map $y=h(q)$ and target $y_*$, the exact task-preserving set is
\begin{equation}
\mathcal M_{y_*}=\{q:h(q)=y_*\}.
\label{eq:ch8:ucm}
\end{equation}
At a regular point $\bar q$ with rank-$r$ Jacobian $J=Dh(\bar q)$, its tangent
space is $\ker J$ and has dimension $n-r$. The local uncontrolled-manifold
analysis uses this linear approximation~\cite{ScholzSchoner1999}. A finite step
in a tangent direction generally leaves the curved level set:
\begin{equation}
h(\bar q+\delta q)-h(\bar q)
=J\delta q+\tfrac12 D^2h(\bar q)[\delta q,\delta q]
+O(\|\delta q\|^3).
\end{equation}
For $h(q)=q_1^2+q_2^2$, $\bar q=(1,0)$ and $\delta q=(0,0.1)$,
$J\delta q=0$ but the task value increases from 1 to 1.01. A singular point
requires a separate analysis; a changing Jacobian rank invalidates one fixed
dimension formula across the whole region.

Using a declared Euclidean metric, let $P_N=I-J^\dagger J$ and
$P_R=J^\dagger J$. For joint-deviation covariance $\Sigma$, dimension-normalized
variances are
\begin{equation}
V_{\mathrm{UCM}}=\frac{\operatorname{tr}(P_N\Sigma)}{n-r},\qquad
V_{\mathrm{ORT}}=\frac{\operatorname{tr}(P_R\Sigma)}{r},
\quad 0<r<n.
\label{eq:ch8:ucm-variance}
\end{equation}
Coordinates mixing lengths and angles require meaningful scaling or another
metric. For isotropic noise $\Sigma=\sigma^2I$, both normalized variances equal
$\sigma^2$; redundancy alone does not give a ratio greater than one. Repeated
trials also need alignment, a task definition and a treatment of measurement
noise. More variance in the tangent directions supports a particular
task-stabilization hypothesis; it does not establish that these directions
are neurologically uncontrolled or irrelevant to every goal.

In golf, a null direction for hand position may change face orientation, shaft
strain, grip loading or a later collision. The chosen output should therefore
reflect the actual question. Kinematic task-preserving variations and muscle
activation synergies live in different spaces; relating them requires muscle
geometry and dynamics rather than identifying the two subspaces by name.

\subsection{Instantaneous Rank Is Not Controllability}

If an input policy imposes $u=Wc$ in $\dot x=f(x)+G(x)u$, the instantaneous
input map is $GW$, with rank at most $k$. Nonnegative and bounded amplitudes
further limit the feasible directions to a cone or bounded set. This does not
mean that the other state directions are uncontrollable over time.

For the double integrator $\dot q=v$, $\dot v=u$, the instantaneous input
matrix is $B=(0,1)^T$, but
\begin{equation}
A=\begin{pmatrix}0&1\\0&0\end{pmatrix},\qquad
[B,AB]=\begin{pmatrix}0&1\\1&0\end{pmatrix}
\label{eq:ch8:controllability-example}
\end{equation}
has rank two. Drift transports the input's effect into position. Along a
general nominal motion, use the linearized pair $A(t),B(t)$ and, where
appropriate, its finite-horizon Gramian. Nonlinear accessibility, finite-time
reachability and controllability under bounds are distinct additional questions.

Nor does a direction invisible in the current output remain unobservable.
For the same double integrator with output $y=q$, velocity is initially in
$\ker C$, $C=(1,0)$, yet $[C^T,(CA)^T]^T$ has rank two. Position measured over
time reveals velocity in the ideal model. There is no theorem that discarded
synergy directions must coincide with an unobservable UCM. The practical test
is whether the restricted policy can meet the chosen task and disturbance
tolerances within its available time and force limits.

\section{Impedance and Co-Contraction}

\subsection{Force, Displacement and Velocity Ports}

For a local scalar perturbation model around a fixed operating point,
\begin{equation}
M\delta\ddot x+D\delta\dot x+K\delta x=\delta F,
\qquad M,D,K>0,
\label{eq:ch8:impedance-time}
\end{equation}
the zero-initial-condition Laplace relations are
\begin{align}
S(s)&=\frac{\delta F(s)}{\delta X(s)}=Ms^2+Ds+K,
\label{eq:ch8:dynamic-stiffness}\\
Z(s)&=\frac{\delta F(s)}{\delta V(s)}=Ms+D+K/s,
\label{eq:ch8:impedance-laplace}\\
H(s)&=\frac{\delta X(s)}{\delta F(s)}=\frac1{Ms^2+Ds+K}.
\end{align}
$S$ is dynamic stiffness (N/m); $Z$ is mechanical impedance (N\,s/m); $H$
is displacement compliance. Some robotics literature uses ``impedance control''
for enforcing the complete mass--damping--stiffness relation. Naming that strategy
does not remove the distinction between the transfer functions and their units.
Multi-axis impedance needs matrices, a frame and a power-consistent port.

The natural frequency is $\omega_n=\sqrt{K/M}$ and damping ratio
$\zeta=D/(2\sqrt{MK})$. Increasing $K$ at fixed $D$ raises $\omega_n$ but lowers
$\zeta$. For the chosen $M=1$ kg, $D=10$ N\,s/m and $K=100$ N/m,
$\omega_n=10$ rad/s and $\zeta=0.5$. Increasing $K$ to 400 N/m doubles the
natural frequency and halves the damping ratio. A faster natural time scale
does not guarantee less overshoot, smaller contact forces or better performance.

\subsection{How Muscle Forces Contribute to Joint Stiffness}

At fixed activation and internal-state assumptions, take
$\tau(q)=-\sum_i F_i(l_i(q))\nabla l_i(q)$. Linearized restoring stiffness is
\begin{equation}
K_q=-\frac{\partial\tau}{\partial q}
=\sum_i\left[
k_i\nabla l_i\nabla l_i^T+F_i\nabla^2l_i\right],
\qquad k_i=\frac{\partial F_i}{\partial l_i}.
\label{eq:ch8:cocontraction-stiffness}
\end{equation}
The first term is material stiffness transmitted through moment arms; the
second is geometric stiffness from loaded, changing geometry. Muscle--tendon
series compliance changes the effective $k_i$. Coupled muscle dynamics,
reflexes and active feedback require a frequency-dependent extension.

In a chosen one-joint model with constant opposing moment arms of 0.03 m,
equal forces of 100 N produce zero net torque. If both effective longitudinal
stiffnesses are 10,000 N/m, the joint stiffness is
$2(0.03)^2(10{,}000)=18$ N\,m/rad. The antagonists carry real load even though
their net torque cancels. Equal activation need not produce equal opposing
torques when muscle sizes, lengths or moment arms differ.

Hogan's analysis and forearm-posture experiment support a role for antagonist
coactivation in impedance modulation~\cite{Hogan1984}. They do not establish
a universal stiffness proportional to target uncertainty. Target uncertainty,
external force disturbances, observation noise and activation-dependent noise
enter different parts of a control problem. Increasing stiffness can resist
a force disturbance while increasing the force caused by a position error
against a stiff environment. The disturbance and the scored outcome must be named.

\subsection{An Explicit Optimization Example}

For a quasistatic scalar model $e=w/K$ with zero-mean force disturbance of
variance $\sigma_F^2$, suppose a designer chooses the surrogate
\begin{equation}
J(K)=\frac{\sigma_F^2}{K^2}+\rho K^2,\qquad
0<K_{\min}\leq K\leq K_{\max},\quad\rho>0.
\label{eq:ch8:impedance-cost}
\end{equation}
Weights carry the units needed to compare the terms. The first term is expected
squared displacement; the second is an assumed stiffness penalty, not measured
metabolic energy. Before applying bounds, differentiation gives
\begin{equation}
K_0=(\sigma_F^2/\rho)^{1/4},\qquad
K^*=\operatorname{clip}(K_0,K_{\min},K_{\max}).
\label{eq:ch8:optimal-k}
\end{equation}
For $\sigma_F=10$ N, $\rho=10^{-6}$ m$^4$/N$^2$ and bounds
20--150 N/m, $K^*=100$ N/m. Four times the force standard deviation would
make the unconstrained optimum 200 N/m and the bounded optimum 150 N/m.
Maximum feasible stiffness can therefore be optimal for a stated task. Changing
the penalty, disturbance spectrum or dynamics changes the result; there is
no general linear ``signal-to-noise'' stiffness law.

For the stable compliance $H$ above,
$\|H\|_\infty=\sup_\omega|H(i\omega)|$ measures a frequency-domain gain and
the induced $L_2$ gain under the usual zero-initial-condition assumptions.
It does not guarantee bounded output for arbitrary unlimited disturbances or
robustness to unspecified model uncertainty. For this scalar system,
\begin{equation}
\|H\|_\infty=
\begin{cases}
1/K,&\zeta\geq1/\sqrt2,\\[2pt]
1/[2\zeta K\sqrt{1-\zeta^2}],&0<\zeta<1/\sqrt2.
\end{cases}
\label{eq:ch8:hinfinity}
\end{equation}
The dynamic-stiffness polynomial $S$ instead grows without bound at high
frequency; it is not the stable proper disturbance-to-displacement transfer
whose finite norm was intended. In a real controller, sensing, filters, delay,
contact and actuator dynamics determine the relevant closed-loop transfer.

\subsection{Mechanical Work and Metabolic Cost}

The local passive model stores
$E=\tfrac12M\delta\dot x^2+\tfrac12K\delta x^2$ and, for constant $M,K$,
obeys $\dot E=\delta F\delta\dot x-D\delta\dot x^2$.
If stiffness is varied, an additional $\tfrac12\dot K\delta x^2$ term enters
the energy balance. An actively changed spring parameter can supply energy;
passivity cannot be inferred from positive instantaneous $K$ alone.

Physiological metabolic cost also depends on activation, contraction history,
force and velocity, tissue properties and the selected metabolic model.
It is not generally proportional to $\int K\,dt$, net joint work, or squared
inverse-dynamics torque. A golfer can have low net wrist torque while opposing
muscles are loaded. A robot can hold a load with zero joint work while its
motor dissipates electrical energy. Report the measured or modeled cost actually
used; mechanical cancellation does not make physiological effort disappear.

\section{Pre-Stress and Stability of Interacting Subsystems}

\subsection{What a Tensegrity Analogy Can Contribute}

Pre-stress means an internal load state present before an additional perturbation.
A self-stress in a free structural model balances internally without external
nodal load. More generally, a loaded structure is linearized around an equilibrium
that includes its external loads. Neither definition says that ``internal forces
exceed external forces''; these are different force collections, not one scalar
comparison.

For a declared elastic potential $U(q)$, the constrained tangent stiffness comes
from the appropriate second variation of total potential at equilibrium. Material
and geometric contributions both matter, as in the preceding muscle-path formula.
Preloading can stiffen some directions and soften others, including buckling.
Equilibrium force balance alone does not prove stability.

Tensegrity models can help formulate questions about distributed load paths.
Anatomical bones also experience bending, shear and joint-contact forces, while
muscles, tendons and fascia have distributed geometry and constitutive behavior.
A cell-scale or ideal strut--cable model is not by itself validation of a
whole-body architecture. Rigid-body multibody models can already include muscle
forces, compliance and loaded constraints; these modeling approaches need not
be treated as mutually exclusive.

For golf, useful questions concern the force path from ground through legs,
trunk, arms and grip, and the resulting tangent response under a defined
perturbation. Muscle tone supplies force and may change that response; it is
not separate from force generation. A proposed fascia-mediated mechanism must
predict an observable effect beyond what the competing joint/contact model
already explains. A humanoid comparison likewise needs a specified structure,
actuation and contact model, rather than an anatomical analogy alone.

\subsection{Why Stable Parts Can Form an Unstable Whole}

Organizing analysis into tissue, intermuscular and task levels is useful.
Their time scales overlap, however, and their couplings must remain in the
equations. A spinal feedback loop is a CNS process, and faster mechanics does
not eliminate sensing, prediction or adaptation. Separate stability does not
guarantee that the sum of the levels' Lyapunov functions is a composite
certificate: coupling terms must also be controlled.

\begin{example}[colframe=chapblue]{An Unstable Interconnection of Stable Scalar Parts}{bio-coupling}
Consider
\begin{equation}
\dot z_1=-z_1+2z_2,\qquad \dot z_2=-z_2+2z_1.
\label{eq:ch8:unstable-coupling}
\end{equation}
Each isolated equation decays when the other state is fixed at zero. The
coupled matrix has eigenvalues $1$ and $-3$. For
$V=(z_1^2+z_2^2)/2$, $\dot V=-z_1^2-z_2^2+4z_1z_2$ is positive at
$z_1=z_2\ne0$. Separate stability has not controlled the feedback interaction.
\end{example}

\begin{theorem}[A Sufficient Composite Decay Condition]
\label{thm:nested-lyapunov}
Suppose positive-definite local functions $V_i$ for a common equilibrium satisfy,
on a specified neighborhood,
\[
\dot V_i\leq-\alpha_i\|z_i\|^2+
\sum_{j\ne i}\gamma_{ij}\|z_i\|\|z_j\|,
\qquad \alpha_i>0,\quad\gamma_{ij}\geq0.
\]
If weights $w_i>0$ make the symmetric matrix with entries
\[
S_{ii}=w_i\alpha_i,\qquad
S_{ij}=-\tfrac12(w_i\gamma_{ij}+w_j\gamma_{ji})
\]
positive definite, then $V=\sum_i w_iV_i$ has strictly negative derivative
away from the equilibrium in that neighborhood. With the usual regularity
and local positive-definiteness bounds, it proves local asymptotic stability.
\end{theorem}
\begin{proof}
Set $r_i=\|z_i\|$. Summing the derivative bounds gives
$\dot V\leq-r^TSr\leq-\lambda_{\min}(S)\|r\|^2$.
Small enough sublevel sets remain inside the neighborhood; the Lyapunov
criterion then applies. A global conclusion would require global hypotheses.
\end{proof}

For three scalar levels with diagonal decay one and symmetric neighbor coupling
0.2, the decay matrix has smallest eigenvalue $1-0.2\sqrt2\simeq0.7172>0$.
This chosen example passes the condition. It is not an estimate of a nervous
system's coupling strength. Other systems may need input-to-state stability,
small-gain, passivity or singular-perturbation analysis. For contact switches,
also check reset-induced changes in the certificate; continuous decay between
events is not a complete hybrid stability proof.

The resulting research claim is specific: suitable coupling and feedback can
make distributed stabilization work. A hierarchy is an analysis strategy,
not automatic protection from dimensionality, delay or instability.

\section{Building and Testing a Coupled Movement Model}

\subsection{Choose the Model for the Question}

Inverse dynamics recovers the net generalized loads consistent with specified
kinematics, inertia and external loads. It does not inherently assume that
muscle activation is instantaneous or that tendons are rigid: those assumptions
belong to a chosen actuator model. Recovering a joint moment does not identify
the individual muscle forces producing it. Ballistic movements can be particularly
sensitive to activation history, elasticity and contact timing; speed alone
does not justify ignoring those states.

A practical model comparison starts with a rigid-body net-load model, then
adds only the states needed to answer the question. A muscle-driven model can
add activation, fiber and tendon dynamics; a humanoid actuator model can add
motor current, transmission compliance and control delay. Flexible shafts and
grips require corresponding deformation states. Both versions should retain
the same task, initial-condition specification and external-load convention.

Check equilibrium, power balance, force bounds, constraint residuals and numerical
convergence before interpreting differences. A more elaborate model can fit
more observations while being less identifiable. Adding springs or feedback
does not guarantee useful attractors, favorable energy transfer or better
predictions. Compare the simplest competing models against held-out data.

\subsection{Identification Requires More Than a Motion Fit}

Let measurements satisfy $y_k=\mathcal H(x(t_k),\theta)+\epsilon_k$, with
parameters $\theta$. The local sensitivity matrix
\begin{equation}
\mathcal S_{kj}=\frac{\partial y_k}{\partial\theta_j}
\label{eq:ch8:identification}
\end{equation}
reveals combinations distinguishable in the selected experiment. When columns
are nearly dependent, different parameters can compensate for one another.
Measurement covariance, parameter scaling and model discrepancy all affect
the interpretation. A successful optimizer does not establish physiological
uniqueness or remove uncertainty.

Joint motion and ground reactions constrain net mechanics. EMG adds timing and
activation-related evidence with its own calibration limitations. Ultrasound
or other tissue measurements can constrain fiber or tendon behavior when the
measurement is appropriate. Perturbation experiments can separate stiffness,
damping and delayed feedback more effectively than repeated unperturbed motion.
No fixed parameter count guarantees identifiability, and imaging is not uniformly
invasive. Specify what each measurement adds and what remains inferred.

For learning and fatigue, parameters or policies may change across trials.
Separate within-trial state estimation from across-trial adaptation. A model
that refits each swing independently can conceal a failure to predict how a
changed grip, tempo, load or support condition changes the response.

\subsection{A Golf and Humanoid Experiment Ladder}

\begin{enumerate}
\item \textbf{Define the delivered state}. Record the clubhead or tool frame,
reference point, contact event and task tolerance. State whether the objective
is face control, speed, balance, disturbance recovery or a combination. Joint
angles alone do not specify all of these outcomes.
\item \textbf{Define comparable initial conditions}. Match or estimate posture,
velocity, activation and elastic state. Distinguish a new command from a
changed starting state. A zero-excitation continuation preserves the existing
activation and strain history; it is not an instruction to make the person
instantaneously relaxed.
\item \textbf{Apply a declared disturbance}. Vary tool inertia, support,
grip loading or a measured external perturbation within an appropriate protocol.
Specify the time and direction. Disturbances applied before and near impact
probe different available correction horizons.
\item \textbf{Measure both task and internal response}. Compare delivered
state, phase error, hand/ground loads and relevant activation or strain evidence.
A successful endpoint with larger opposing forces may trade robustness for
effort. Lower joint torque alone cannot decide whether that trade occurred.
\item \textbf{Compare mechanisms}. Test a net-load model, an augmented tissue
or actuator model, and alternative feedback policies using the same scored
data. Report prediction intervals and failures, including cases in which added
compliance or reduced input dimension worsens the task.
\end{enumerate}

These comparisons expose how the components are intertwined. Posture changes
moment arms and inertia; activation and strain change force and local response;
contact changes feasible reactions; those changes alter the useful feedback
directions and remaining time. Earlier work can become later kinetic or elastic
energy, while a change that improves speed may impair face control or loading.
The joint explanation is a sequence of measurable consequences, not a claim
that mechanics replaces neural control.

\subsection{What Would Count Against the Proposed Explanation?}

A preflex-based explanation weakens if the estimated intrinsic response is
too small or has the wrong sign to explain early recovery. A tendon-recoil
explanation weakens if the bounded recoverable strain energy is insufficient
or is released at the wrong time. A fixed-synergy explanation weakens if it
fails to predict withheld task directions or perturbations that an independently
actuated model predicts. An attractor explanation weakens if the supposed
recovery disappears after phase alignment or depends on unmodeled powered
feedback. A hierarchical stability claim fails if cross-coupling or contact
resets defeat its proposed certificate.

These are model-specific tests, not reasons to reject biological coordination
as a whole. They make the useful positive argument stronger: a successful
model explains both why a pattern works in its operating region and why it
stops working when the relevant conditions change.

\section{Connection to the Series}

Geometry supplies configurations, tangent spaces and moment arms. Multibody
dynamics translates applied loads into motion. Muscle and actuator models
describe how command history becomes load and impedance. Feedback, prediction
and learning shape behavior within those physical limits. Contact and impact
determine how the delivered state becomes task outcome. Each layer adds an
assumption and a corresponding opportunity for validation.

The neighboring treatments of passive control, central pattern generation,
motor learning and robotics should be read with this separation in mind.
Distributed mechanics and neural computation can cooperate; a low-dimensional
description can be useful without uniquely identifying a neural implementation.
The strongest golf-swing argument specifies the state, intervention, constraints,
time horizon and observable consequence. The same discipline supports a
humanoid-tool model while preserving the differences in its actuators and sensing.

\section*{Exercises}

\begin{enumerate}
\item \textbf{Force--Velocity and Series Equilibrium.}
Verify the values and one-sided derivatives of the stated $f_V$ at zero.
Use the declared force--length and tension-only tendon laws with zero pennation
to solve a prescribed length-history example. Increasing fiber length is
lengthening, not shortening. Report how tendon compliance and activation time
scale separately affect the result; do not prescribe inconsistent fiber and
tendon lengths independently.

\item \textbf{Stored Energy and Zero Excitation.}
Reproduce the 4 J tendon example and its ideal average recoil power. Then use
$a(t_0)=0.8$ and a chosen $\tau_a=0.04$ s to calculate activation 0.04 s after
zeroing excitation. Explain why the resulting $0.8/e\simeq0.2943$ need not
produce zero active force. Close the energy balance when losses are added.

\item \textbf{Orbit and Trajectory Stability.}
Simulate the radial/phase example from two nearby initial phases and radii.
Plot distance to the orbit and time-matched distance between trajectories.
Explain which measure is appropriate for a periodic rhythm and which additional
timing information a collision task requires.

\item \textbf{HKB Transition.}
Use the first-order phase equation with $a_H=1$ and decrease $b_H$ through
0.25. Start near, rather than exactly at, $\pi$. Compare the observed recovery
rate with $a_H-4b_H$ and test sensitivity to perturbation size and parameter-ramp
speed. A transition frequency cannot be reported until a frequency-to-parameter
relationship has been specified.

\item \textbf{Uncontrolled-Manifold Null Model.}
At a regular configuration of a three-link planar arm, compute a rank-two
position Jacobian and the dimension-normalized covariance projections.
First use isotropic joint noise: the population ratio is one, not automatically
greater than one. Then impose a larger covariance in the tangent direction.
Compare linearized task variation with the actual nonlinear forward kinematics
as perturbation size increases. Repeat with orientation included in the task.

\item \textbf{Impedance and Coupling Trade-Offs.}
Reproduce the bounded stiffness optimum, then compare disturbance compliance
and contact force under a prescribed position error. Verify the two-level
unstable interconnection and the three-level sufficient decay condition.
Finally introduce a delay or reset and determine which earlier certificate
must be reconsidered. A favorable result from one objective is not a universal
ranking of stiff and compliant strategies.
\end{enumerate}

\chapter{Central Pattern Generators}
\label{ch:cpg}

\epigraph{%
  Rhythm is the most ancient form of motor control.
}{}

\section{CPGs in Biological Motor Control}

Central pattern generators are neural circuits that produce rhythmic motor output
without rhythmic sensory input or supraspinal commands. They underlie locomotion,
respiration, chewing, and other rhythmic behaviors.

\section{The Matsuoka Oscillator}

The most widely used CPG model in robotics is the Matsuoka oscillator
\cite{Matsuoka1985}:
\begin{align}
  \tau_1 \dot{x}_i &= -x_i - \beta v_i - \sum_{j \neq i} w_{ij} y_j + u_i + s_i, \\
  \tau_2 \dot{v}_i &= -v_i + y_i, \\
  y_i &= \max(0, x_i),
  \label{eq:ch8:matsuoka}
\end{align}
where $x_i$ is the membrane potential, $v_i$ is the adaptation variable, $y_i$ is the
output, $w_{ij}$ are inhibitory coupling weights, $u_i$ is the tonic input, and
$s_i$ is the sensory feedback.

\begin{lstlisting}[language=Python, caption={Matsuoka oscillator CPG implementation.}]
import numpy as np
from scipy.integrate import solve_ivp

class MatsuokaOscillator:
    """Two-neuron Matsuoka oscillator for CPG-based control."""

    def __init__(
        self,
        tau1: float = 0.1,
        tau2: float = 0.3,
        beta: float = 2.5,
        w_mutual: float = 2.0,
        tonic_input: float = 1.0
    ):
        self.tau1 = tau1
        self.tau2 = tau2
        self.beta = beta
        self.w = w_mutual
        self.u = tonic_input

    def dynamics(self, t: float, state: np.ndarray) -> np.ndarray:
        x1, v1, x2, v2 = state
        y1 = max(0.0, x1)
        y2 = max(0.0, x2)

        dx1 = (-x1 - self.beta * v1 - self.w * y2 + self.u) / self.tau1
        dv1 = (-v1 + y1) / self.tau2
        dx2 = (-x2 - self.beta * v2 - self.w * y1 + self.u) / self.tau1
        dv2 = (-v2 + y2) / self.tau2

        return np.array([dx1, dv1, dx2, dv2])

    def simulate(self, t_end: float = 5.0,
                 dt: float = 0.001) -> tuple[np.ndarray, np.ndarray]:
        state0 = np.array([0.1, 0.0, -0.1, 0.0])
        t_eval = np.arange(0, t_end, dt)
        sol = solve_ivp(self.dynamics, [0, t_end], state0,
                        t_eval=t_eval, method='RK45')

        # Extract outputs
        y1 = np.maximum(0, sol.y[0])
        y2 = np.maximum(0, sol.y[2])
        output = y1 - y2  # Net output

        return sol.t, output


# Simulate and print summary
cpg = MatsuokaOscillator(tau1=0.1, tau2=0.3)
t, output = cpg.simulate(t_end=5.0)

# Estimate frequency from zero crossings
crossings = np.where(np.diff(np.sign(output)))[0]
if len(crossings) >= 4:
    periods = np.diff(t[crossings[::2]])
    freq = 1.0 / np.mean(periods)
    print(f"CPG frequency: {freq:.2f} Hz")
    print(f"CPG amplitude: {np.max(np.abs(output[len(output)//2:])):.3f}")
\end{lstlisting}

\section{Phase Coupling and Inter-Limb Coordination}

Multiple CPGs can be coupled to produce coordinated multi-limb movement:
\begin{equation}
  \dot{\phi}_i = \omega_i + \sum_{j} K_{ij} \sin(\phi_j - \phi_i - \psi_{ij}),
  \label{eq:ch8:phase-coupling}
\end{equation}
where $\phi_i$ is the phase of the $i$-th oscillator, $\omega_i$ is its natural
frequency, $K_{ij}$ is the coupling strength, and $\psi_{ij}$ is the desired phase
relationship.

For quadruped locomotion:
\begin{itemize}
  \item Walk: $\psi = (0, \pi, \pi/2, 3\pi/2)$ (diagonal pairs alternating)
  \item Trot: $\psi = (0, \pi, \pi, 0)$ (diagonal pairs synchronized)
  \item Gallop: $\psi = (0, 0, \pi, \pi)$ (front-back alternating)
\end{itemize}

\section*{Exercises}

\begin{enumerate}
  \item \textbf{CPG Tuning.} Vary $\tau_1/\tau_2$ ratio and plot the resulting
        frequency and duty cycle. How does the ratio control the waveform shape?

  \item \textbf{Phase Coupling.} Implement 4 coupled oscillators and produce
        walk, trot, and gallop patterns by varying $\psi_{ij}$.

  \item \textbf{(Programming.)} Use the Matsuoka CPG to drive a simulated bipedal
        walker. Tune the CPG parameters for stable walking.
\end{enumerate}

\chapter{Motor Learning and Adaptation}
\label{ch:motor-learning}

\epigraph{%
  Practice does not make perfect.
  Practice makes permanent.
  Only perfect practice makes perfect.
}{}

\section{Three Learning Systems}

Motor learning involves three distinct neural systems with different timescales and mechanisms:

\begin{enumerate}
  \item \textbf{Error-based learning} (cerebellum): uses sensory prediction errors (climbing
        fiber signals) to update cerebellar forward models via LTD (Marr, 1969; Kawato et al.,
        1987). Trial-by-trial timescale (minutes to hours). Dominant mechanism for visuomotor
        adaptation (Shadmehr and Mussa-Ivaldi, 1994).

  \item \textbf{Reinforcement learning} (basal ganglia): uses dopamine reward signals to
        strengthen successful motor patterns (Schultz et al., 1997). Slower than error-based
        learning (hours to days). Enables discovery of novel solutions not implicit in the
        task structure.

  \item \textbf{Consolidation and structural learning} (cortical): repetition of a movement
        biases future movements toward repeated patterns (Cohen et al., 2005). Timescale of
        hours to days; explains habit formation. Involves gradual changes in M1 tuning and
        synaptic weights.
\end{enumerate}

\begin{figure}[h]
\centering
\begin{tikzpicture}[scale=0.85]
  % Timeline
  \draw (0, 0) -- (8, 0);
  \draw (0, -0.2) node[align=center, anchor=north] {Baseline};
  \draw (3.5, -0.2) node[align=center, anchor=north] {Adaptation};
  \draw (7.5, -0.2) node[align=center, anchor=north] {Washout};

  % Error curve (schematic)
  \draw[blue, thick] (0, 1) -- (1.5, 1) -- (2, 3) -- (3.5, 1.2) -- (4.5, 0.3) -- (5.5, 0.1) -- (8, 0.8);

  % Labels
  \draw (0.75, 0.5) node[align=center, anchor=center, font=\tiny] {minimal};
  \draw (2.75, 3.3) node[align=center, anchor=center, font=\tiny] {early};
  \draw (4, 0.8) node[align=center, anchor=center, font=\tiny] {adaptation};
  \draw (7.2, 1.2) node[align=center, anchor=center, font=\tiny] {aftereffect};

  % Dashed line (zero)
  \draw[dashed] (0, 0) -- (8, 0);

  % Y-axis
  \draw[->, thick] (-0.3, 0) -- (-0.3, 3.5) node[align=center, anchor=east, rotate=90] {Error};

  % Title
  \draw (4, 4.2) node[align=center, anchor=center, font=\small\bf] {Visuomotor Adaptation Curve};
  \draw (4, 3.8) node[align=center, anchor=center, font=\tiny] {(Reaching with visual rotation)};
\end{tikzpicture}
\caption{Characteristic visuomotor adaptation profile (Shadmehr and Mussa-Ivaldi, 1994):
rapid error reduction (first 20 trials), gradual asymptote, persistent aftereffect
demonstrating learning consolidation.}
\label{fig:ch9:visuomotor-adaptation}
\end{figure}

\section{Visuomotor Adaptation}

The classic paradigm for studying error-based learning (Kamin, 1969; Shadmehr and Mussa-Ivaldi,
1994):

\begin{lstlisting}[language=Python, caption={Visuomotor adaptation simulation.}]
import numpy as np

def simulate_visuomotor_adaptation(
    n_baseline: int = 50,
    n_adaptation: int = 200,
    n_washout: int = 100,
    perturbation_deg: float = 30.0,
    learning_rate: float = 0.05,
    retention: float = 0.98
) -> dict:
    """Simulate visuomotor rotation adaptation.

    A simple state-space model of adaptation:
    x_{k+1} = retention * x_k + learning_rate * e_k

    Parameters
    ----------
    n_baseline : int
        Baseline trials (no perturbation).
    n_adaptation : int
        Adaptation trials (perturbation on).
    n_washout : int
        Post-adaptation trials (perturbation off).
    perturbation_deg : float
        Rotation perturbation in degrees.
    learning_rate : float
        Rate of error-based learning.
    retention : float
        Memory retention between trials.

    Returns
    -------
    dict
        Keys: 'trial', 'reaching_error', 'internal_state',
              'phase' (labels for each trial)
    """
    n_total = n_baseline + n_adaptation + n_washout
    p = np.radians(perturbation_deg)

    # Internal state (learned compensation)
    x = 0.0
    errors = np.zeros(n_total)
    states = np.zeros(n_total)
    phases = []

    for trial in range(n_total):
        if trial < n_baseline:
            perturbation = 0.0
            phases.append('baseline')
        elif trial < n_baseline + n_adaptation:
            perturbation = p
            phases.append('adaptation')
        else:
            perturbation = 0.0
            phases.append('washout')

        # Reaching error = perturbation - compensation
        error = perturbation - x
        errors[trial] = np.degrees(error)
        states[trial] = np.degrees(x)

        # Update internal state
        x = retention * x + learning_rate * error

    return {
        'trial': np.arange(n_total),
        'reaching_error': errors,
        'internal_state': states,
        'phase': phases
    }


result = simulate_visuomotor_adaptation()
print("Visuomotor Adaptation Summary:")
print(f"  Baseline error:  {np.mean(np.abs(result['reaching_error'][:50])):.1f} deg")
print(f"  Final adaptation error: "
      f"{np.abs(result['reaching_error'][249]):.1f} deg")
print(f"  Aftereffect: {result['reaching_error'][250]:.1f} deg")
\end{lstlisting}

\section{Savings and Interference}

\subsection{Savings}
Re-learning a previously learned adaptation is faster than initial learning,
even after complete washout (Ebbinghaus, 1885; Huang et al., 2011). This savings
effect suggests that motor memories are not erased but suppressed, with additional
learning circuits retaining implicit information about the learning dynamics.

\subsection{Interference and Dual-Rate Learning}
Learning opposite perturbations (A then B) produces anterograde interference:
learning B slows retention of A (Li and Xie, 2000). This is explained by dual-rate
learning models: a fast process learns trial-by-trial; a slow process consolidates.
Interference occurs because the fast process is task-specific while the slow process
is more general.

\section{Structural Learning and Meta-Learning}

Over many adaptation experiences, subjects develop the ability to adapt faster---they
learn to learn (context-dependent learning). Smith et al. (1998) showed that prior
exposure to one visuomotor rotation task accelerates learning of a novel rotation.
This involves learning the structure of the task space itself, not just individual
parameter values. Brain areas like dorsolateral prefrontal cortex appear to encode
such meta-learning (Gershman et al., 2017).

\section*{Exercises}

\begin{enumerate}
  \item \textbf{State-Space Model.} Using the provided code, explore how learning
        rate and retention affect adaptation speed and aftereffect size.

  \item \textbf{Dual-Rate Model.} Implement a dual-rate model (fast + slow process)
        and show that it can explain savings.

  \item \textbf{(Programming.)} Simulate interference by adapting to $+30^\circ$ then
        $-30^\circ$. Plot the adaptation curves and aftereffects.
\end{enumerate}

\chapter{Computational Models of Motor Control}
\label{ch:computational}

\epigraph{%
  A theory of motor control must ultimately be a theory of
  computation: what is computed, how, and why.
}{}

\section{Optimal Feedback Control}

The dominant computational theory of motor control is optimal feedback control (OFC)
\cite{TodorovJordan2002}:

\begin{definition}{The OFC Framework}{ofc-def}
  The nervous system controls movement by solving, in real-time, the stochastic
  optimal control problem:
  \begin{equation}
    \min_{\control(\cdot)} \mathbb{E}\left[
    \phi(\state_T) + \int_0^T \ell(\state_t, \control_t) \dt
    \right],
    \label{eq:ch10:ofc}
  \end{equation}
  subject to:
  \begin{equation}
    \dd\state = f(\state, \control) \dt + G(\state) \dd\bm{w},
  \end{equation}
  where $\dd\bm{w}$ is Brownian noise representing signal-dependent motor noise.
\end{definition}

\noindent
This is precisely the optimal control problem studied in Agrachev \& Sachkov \cite{AgrachevSachkov2004},
Chapters~5--6 (stochastic optimal control). Their framework establishes existence and uniqueness
of optimal feedback laws, characterizes them via the Hamilton-Jacobi-Bellman equation, and shows
that feedback synthesis via optimal control naturally leads to closed-loop stability.

\subsection{The Minimum Intervention Principle}

OFC naturally explains why the nervous system tolerates task-irrelevant variability
(the UCM from Chapter~1): the optimal controller only corrects deviations that
affect the cost function. Task-irrelevant deviations are ``free'' and need not be
corrected.

\begin{lstlisting}[language=Python, caption={Linear-quadratic optimal feedback controller.}]
import numpy as np
from scipy.linalg import solve_continuous_are

def lqr_ofc(
    A: np.ndarray,
    B: np.ndarray,
    Q: np.ndarray,
    R: np.ndarray,
    state_noise_cov: np.ndarray | None = None
) -> tuple[np.ndarray, np.ndarray]:
    """Compute LQR optimal feedback gains.

    This implements the infinite-horizon OFC solution.

    Parameters
    ----------
    A : np.ndarray, shape (n, n)
        State dynamics matrix.
    B : np.ndarray, shape (n, m)
        Input matrix.
    Q : np.ndarray, shape (n, n)
        State cost matrix.
    R : np.ndarray, shape (m, m)
        Control cost matrix.
    state_noise_cov : np.ndarray, shape (n, n), optional
        Process noise covariance.

    Returns
    -------
    tuple
        (K, P) where K is the optimal gain and P is the Riccati solution.
    """
    # Solve continuous algebraic Riccati equation
    P = solve_continuous_are(A, B, Q, R)

    # Optimal gain
    K = np.linalg.solve(R, B.T @ P)

    return K, P


# Example: 2D reaching with motor noise
# State: [x, y, dx, dy], Control: [Fx, Fy]
A = np.array([
    [0, 0, 1, 0],
    [0, 0, 0, 1],
    [0, 0, -2, 0],  # damping
    [0, 0, 0, -2],
])
B = np.array([
    [0, 0],
    [0, 0],
    [1, 0],
    [0, 1],
])

# Task: reach to target with minimum effort
# High cost on position at target, low cost on path
Q = np.diag([10, 10, 1, 1])  # Position matters more than velocity
R = 0.01 * np.eye(2)  # Low control cost

K, P = lqr_ofc(A, B, Q, R)
print("OFC gain matrix K:")
print(K)
print("\nMinimum intervention: gains are higher for")
print("task-relevant (position) than task-irrelevant dimensions")
\end{lstlisting}

\section{Signal-Dependent Noise}

A key feature of biological motor systems is that motor noise is
\textbf{signal-dependent}: the noise variance scales with the control signal magnitude:
\begin{equation}
  \control = \bar{\control} + \sigma(\bar{\control}) \bm{\xi},
  \quad \sigma(\bar{\control}) = k \norm{\bar{\control}},
  \label{eq:ch10:signal-dep-noise}
\end{equation}
where $\bm{\xi}$ is white noise and $k$ is the noise coefficient.

This explains minimum-jerk trajectories: the optimal trajectory under signal-dependent
noise minimizes the time-integral of squared jerk \cite{FlashHogan1985, Harris1998}.

\section{Active Inference (Friston Framework)}

Active inference reframes motor control as inference: the agent infers the
actions that would make its predictions come true:
\begin{equation}
  F = \mathbb{E}_q[\ln q(\state) - \ln p(\state, \bm{y})],
  \label{eq:ch10:free-energy}
\end{equation}
where $F$ is the variational free energy, $q(\state)$ is the approximate posterior,
and $p(\state, \bm{y})$ is the generative model.

\section*{Exercises}

\begin{enumerate}
  \item \textbf{OFC Reaching.} Using the LQR code, simulate reaching movements
        to different targets. Show that task-irrelevant variability is tolerated.

  \item \textbf{Signal-Dependent Noise.} Simulate reaching with and without
        signal-dependent noise. Plot the trajectory variability as a function of
        movement speed.

  \item \textbf{(Programming.)} Implement the minimum-jerk model and the
        minimum-variance model. Show that both predict smooth bell-shaped velocity
        profiles.
\end{enumerate}

\chapter{From Neural Architecture to Robot Control}
\label{ch:neural-to-robot}

\epigraph{%
  If we truly understood how humans move,
  we could build robots that move like humans.
  We cannot. This tells us something.
}{}

\section{Lessons from Biology for Robot Design}

The four volumes of theory converge on practical design principles for robotic
control systems:

\begin{insight}{Bio-Inspired Control Design Principles}{bio-robot-design}
  \begin{enumerate}
    \item \textbf{Exploit passive dynamics} (Vol~II, Vol~IV Ch~7): design the mechanical
          system so that natural dynamics do useful work. Use compliance, not rigidity.
    \item \textbf{Use shallow-wide controllers} (Vol~IV Ch~4): for real-time control,
          prefer 2--3 layer networks with thousands of neurons over deep networks with
          hundreds of layers.
    \item \textbf{Employ synergy-based dimensionality reduction} (Vol~IV Ch~4): define
          4--8 actuation synergies and control in synergy space, not muscle space.
    \item \textbf{Implement forward models} (Vol~IV Ch~6): use the Smith predictor
          architecture to compensate for sensor/actuator delays.
    \item \textbf{Design for contraction} (Vol~I Ch~4): ensure the closed-loop
          system is contracting, guaranteeing stability and robustness.
    \item \textbf{Optimize trajectories, not setpoints} (Vol~II): design
          trajectory-centric controllers with funnel guarantees.
    \item \textbf{Accept variability} (Vol~IV Ch~1): allow task-irrelevant
          variability. Don't over-constrain the controller.
  \end{enumerate}
\end{insight}

\section{Cerebellar-Inspired Adaptive Control}

The cerebellum's architecture suggests a specific adaptive controller:
\begin{itemize}
  \item \textbf{Granule cells} $\to$ basis function expansion (like kernel methods)
  \item \textbf{Purkinje cells} $\to$ linear combination of basis functions
  \item \textbf{Climbing fiber error} $\to$ supervised learning signal
\end{itemize}

This maps directly to:
\begin{equation}
  \control_{\text{adaptive}}(\state) = \hat{W}^T \bm{\phi}(\state),
  \quad \dot{\hat{W}} = -\Gamma \bm{\phi}(\state) \bm{e}^T,
  \label{eq:ch11:cerebellar-control}
\end{equation}
where $\bm{\phi}(\state)$ are radial basis functions (granule cell expansion),
$\hat{W}$ are adaptive weights (parallel fiber--Purkinje cell synapses), and
$\bm{e}$ is the tracking error (climbing fiber signal).

\section{CPG-Based Locomotion Controllers}

Following Chapter~8, CPG-based locomotion controllers use coupled oscillators to
generate walking/running patterns. The CPG provides the \emph{timing} and the
musculoskeletal model provides the \emph{mechanics}:

\begin{lstlisting}[language=Python, caption={Bio-inspired locomotion controller architecture.}]
import numpy as np

class BioInspiredLocomotionController:
    """Hierarchical control architecture inspired by
    biological motor control."""

    def __init__(
        self,
        n_joints: int,
        n_synergies: int = 4,
        cpg_freq: float = 1.0
    ):
        self.n_joints = n_joints
        self.n_synergies = n_synergies
        self.cpg_freq = cpg_freq

        # Synergy matrix (learned from motion data)
        self.W = np.random.randn(n_joints, n_synergies) * 0.1

        # CPG phases for each synergy
        self.phases = np.linspace(
            0, 2 * np.pi * (1 - 1/n_synergies), n_synergies
        )

        # Feedback gains (shallow: single layer)
        self.K_fb = np.zeros((n_synergies, 2 * n_joints))

    def cpg_output(self, t: float) -> np.ndarray:
        """Generate rhythmic synergy activations from CPG."""
        activations = np.array([
            max(0.0, np.sin(2 * np.pi * self.cpg_freq * t + phi))
            for phi in self.phases
        ])
        return activations

    def feedforward(self, t: float) -> np.ndarray:
        """Feedforward joint torques from CPG + synergies."""
        c = self.cpg_output(t)
        return self.W @ c

    def feedback(self, state_error: np.ndarray) -> np.ndarray:
        """Feedback correction (shallow: one layer)."""
        c_fb = self.K_fb @ state_error
        return self.W @ c_fb

    def control(self, t: float,
                state_error: np.ndarray) -> np.ndarray:
        """Total control: feedforward (CPG) + feedback."""
        return self.feedforward(t) + self.feedback(state_error)


# Example: 6-joint biped with 4 synergies
controller = BioInspiredLocomotionController(
    n_joints=6, n_synergies=4, cpg_freq=1.2
)

# Generate one cycle of control signals
t = np.linspace(0, 1.0/1.2, 200)
torques = np.array([controller.feedforward(ti) for ti in t])

print(f"Bio-inspired controller: {controller.n_joints} joints, "
      f"{controller.n_synergies} synergies")
print(f"Control dimensionality reduced from {controller.n_joints} "
      f"to {controller.n_synergies}")
print(f"Peak torque magnitude: {np.max(np.abs(torques)):.3f}")
\end{lstlisting}

\section{Whole-Body Control Architectures}

For humanoid robots, the bio-inspired architecture suggests:

\begin{enumerate}
  \item \textbf{Level 0 (Spinal)}: joint-level impedance control with reflex-like
        feedback ($\sim$1~ms loop)
  \item \textbf{Level 1 (CPG)}: rhythmic pattern generation for locomotion
        ($\sim$10~ms loop)
  \item \textbf{Level 2 (Cerebellar)}: adaptive feedforward based on forward model
        predictions ($\sim$50~ms loop)
  \item \textbf{Level 3 (Cortical)}: trajectory planning and task-level control
        ($\sim$100~ms loop)
\end{enumerate}

Each level is \textbf{shallow} (1--2 layers) but \textbf{wide} (many parallel channels),
operating at progressively longer timescales.

\section{The Connection to UpstreamDrift}

Volume~V provides the practical tools to implement and test these controllers using
the UpstreamDrift simulation platform. The connection points:
\begin{itemize}
  \item \textbf{MuJoCo}: test CPG-based locomotion controllers on musculoskeletal models
  \item \textbf{Drake}: implement trajectory optimization for trajectory-centric control
  \item \textbf{Pinocchio}: compute forward/inverse dynamics efficiently in real-time
\end{itemize}

\section*{Exercises}

\begin{enumerate}
  \item \textbf{Controller Architecture.}
        Using the bio-inspired controller code, vary the number of synergies and
        measure the resulting control signal dimensionality and smoothness.

  \item \textbf{Cerebellar Adaptive Control.}
        Implement the cerebellar-inspired adaptive controller for a 2-DOF arm and
        show that it learns an inverse model over 50 trials.

  \item \textbf{Hierarchical Design.}
        Design a 3-level hierarchical controller for a simulated biped, with
        impedance control at Level 0, CPG at Level 1, and trajectory correction at Level 2.

  \item \textbf{(Programming.)} Implement the full bio-inspired controller for a
        MuJoCo walking model and benchmark against a standard PD controller.
\end{enumerate}


\backmatter
\printindex

\bibliographystyle{plain}
\bibliography{../geometry_of_motion}

\end{document}
